% =====================================================================
% DOCUMENTAÇÃO TÉCNICA E ARQUITETURAL — PORTAL PROGWEB 2
% Estilo visual institucional: preambulo-roteiro-moderno.tex
% Compilação: tectonic arquitetura-modular-frontend.tex
% =====================================================================
\newcommand{\numeroRoteiro}{DOC}
\newcommand{\tituloCurto}{Arquitetura Modular \& Estrutura do Portal}
% =====================================================================
% PREÂMBULO PADRÃO — ROTEIROS PRÁTICOS DE PROGWEB 2 (visual moderno)
% UC Programação para a Internet 2 — IFSC Câmpus Garopaba
% Adaptado do padrão da UC Iniciação Científica e Extensionista
% Documentação: base-de-conhecimento/padrao-visual-roteiros-pdf.md
%
% Uso no main.tex de cada roteiro (o \documentclass está aqui dentro):
%   \newcommand{\numeroRoteiro}{05}
%   \newcommand{\tituloCurto}{Fetch API \& Flask}
%   % =====================================================================
% PREÂMBULO PADRÃO — ROTEIROS PRÁTICOS DE PROGWEB 2 (visual moderno)
% UC Programação para a Internet 2 — IFSC Câmpus Garopaba
% Adaptado do padrão da UC Iniciação Científica e Extensionista
% Documentação: base-de-conhecimento/padrao-visual-roteiros-pdf.md
%
% Uso no main.tex de cada roteiro (o \documentclass está aqui dentro):
%   \newcommand{\numeroRoteiro}{05}
%   \newcommand{\tituloCurto}{Fetch API \& Flask}
%   % =====================================================================
% PREÂMBULO PADRÃO — ROTEIROS PRÁTICOS DE PROGWEB 2 (visual moderno)
% UC Programação para a Internet 2 — IFSC Câmpus Garopaba
% Adaptado do padrão da UC Iniciação Científica e Extensionista
% Documentação: base-de-conhecimento/padrao-visual-roteiros-pdf.md
%
% Uso no main.tex de cada roteiro (o \documentclass está aqui dentro):
%   \newcommand{\numeroRoteiro}{05}
%   \newcommand{\tituloCurto}{Fetch API \& Flask}
%   \input{../base-de-conhecimento/preambulo-roteiro-moderno.tex}
%   \begin{document}
%   \input{content.tex}
%   \end{document}
% Compilar com: tectonic main.tex
% =====================================================================
\documentclass[11pt,a4paper]{article}

\providecommand{\numeroRoteiro}{0X}
\providecommand{\tituloCurto}{Título curto do roteiro}

% =====================================================
% PACOTES
% =====================================================
\usepackage[utf8]{inputenc}
\usepackage[T1]{fontenc}
\usepackage[brazil]{babel}
\shorthandoff{"}
\AtBeginDocument{\shorthandoff{"}}
\usepackage[default]{sourcesanspro}          % texto: sem serifa, leitura confortável
\usepackage[scale=0.86]{sourcecodepro}       % código: monoespaçada moderna
% =====================================================
% ÍCONES (fontawesome5 com compatibilidade Tectonic no macOS)
% Evita crash de alocação na rotina nativa \XeTeXglyphname do Tectonic
% =====================================================
\usepackage{expl3}
\ExplSyntaxOn
\cs_set:Npn \tex_XeTeXglyphname:D #1#2 { }
\ExplSyntaxOff
\newcommand{\faDice}{}
\usepackage{fontawesome5}

\NewDocumentCommand\definirFaIcone{m m}{%
  \expandafter\NewDocumentCommand\csname fa#1\endcsname{O{}}{%
    \if\relax\detokenize{##1}\relax
      \faIcon{#2}%
    \else
      \faIcon[##1]{#2}%
    \fi
  }%
}
\definirFaIcone{AngleRight}{angle-right}
\definirFaIcone{BookOpen}{book-open}
\definirFaIcone{Bullseye}{bullseye}
\definirFaIcone{Calculator}{calculator}
\definirFaIcone{ChalkboardTeacher}{chalkboard-teacher}
\definirFaIcone{CheckCircle}{check-circle}
\definirFaIcone{ChevronRight}{chevron-right}
\definirFaIcone{Circle}{circle}
\definirFaIcone{ClipboardCheck}{clipboard-check}
\definirFaIcone{Code}{code}
\definirFaIcone{Columns}{columns}
\definirFaIcone{Database}{database}
\definirFaIcone{DraftingCompass}{drafting-compass}
\definirFaIcone{ExclamationTriangle}{exclamation-triangle}
\definirFaIcone{ExternalLinkAlt}{external-link-alt}
\definirFaIcone{Flag}{flag}
\definirFaIcone{FolderOpen}{folder-open}
\definirFaIcone{Github}{github}
\definirFaIcone{LaptopCode}{laptop-code}
\definirFaIcone{LayerGroup}{layer-group}
\definirFaIcone{Lightbulb}{lightbulb}
\definirFaIcone{Link}{link}
\definirFaIcone{Lock}{lock}
\definirFaIcone{Map}{map}
\definirFaIcone{MapMarkerAlt}{map-marker-alt}
\definirFaIcone{MousePointer}{mouse-pointer}
\definirFaIcone{Server}{server}
\definirFaIcone{Square}{square}
\definirFaIcone{Terminal}{terminal}
\definirFaIcone{Tools}{tools}
\definirFaIcone{University}{university}
\definirFaIcone{Check}{check}
\definirFaIcone{Times}{times}
\definirFaIcone{InfoCircle}{info-circle}
\definirFaIcone{QuestionCircle}{question-circle}
\definirFaIcone{Play}{play}
\definirFaIcone{FileCode}{file-code}
\definirFaIcone{FileAlt}{file-alt}
\definirFaIcone{Folder}{folder}
\definirFaIcone{Cogs}{cogs}
\definirFaIcone{Cog}{cog}
\definirFaIcone{Eye}{eye}
\definirFaIcone{Desktop}{desktop}
\definirFaIcone{Globe}{globe}
\definirFaIcone{Search}{search}
\definirFaIcone{Download}{download}
\definirFaIcone{Upload}{upload}
\definirFaIcone{Sync}{sync}
\definirFaIcone{ShareAlt}{share-alt}
\definirFaIcone{Edit}{edit}
\definirFaIcone{Trash}{trash}
\definirFaIcone{Plus}{plus}
\definirFaIcone{Minus}{minus}
\usepackage{geometry}
\usepackage{graphicx}
\graphicspath{{../figures/}{figures/}}
\usepackage{fancyhdr}
\usepackage{titlesec}
\usepackage[table]{xcolor}
\usepackage{hyperref}
\usepackage{xurl}                            % URLs longas quebram em qualquer ponto
\usepackage{enumitem}
\usepackage[most]{tcolorbox}
\usepackage{float}
\usepackage{tikz}
\usetikzlibrary{arrows.meta, positioning, shapes.geometric, calc, babel}
\usepackage{tabularx}
\usepackage{booktabs}
\usepackage{listings}
\usepackage{upquote}
\usepackage{lastpage}
\usepackage{qrcode}
\usepackage{multicol}
\usepackage{setspace}
\usepackage{needspace}
\usepackage{adjustbox}

\geometry{top=2.6cm, bottom=2.6cm, left=2cm, right=2cm, headsep=0.8cm}
\linespread{1.18}
\setlength{\parskip}{0.55em}
\setlength{\parindent}{0pt}
\setlength{\emergencystretch}{6em}   % evita linhas estouradas com \texttt longos

% =====================================================
% PALETA (cores institucionais IFSC, em tons suaves)
% =====================================================
\definecolor{azul}{RGB}{30,60,90}
\definecolor{azulMedio}{RGB}{37,99,235}
\definecolor{verde}{RGB}{22,101,52}
\definecolor{verdeVivo}{RGB}{22,163,74}
\definecolor{ambar}{RGB}{217,119,6}
\definecolor{vermelho}{RGB}{220,38,38}
\definecolor{roxo}{RGB}{124,58,237}
\definecolor{laranja}{RGB}{194,65,12}
\definecolor{ciano}{RGB}{14,116,144}
\definecolor{tinta}{RGB}{30,41,59}
\definecolor{suave}{RGB}{100,116,139}
\definecolor{borda}{RGB}{226,232,240}
\definecolor{fundo}{RGB}{248,250,252}
\definecolor{barraTitulo}{RGB}{241,245,249}
\definecolor{terminal}{RGB}{15,23,42}

% Nomes antigos dos roteiros de ProgWeb 2 (mantidos para compatibilidade)
\colorlet{azulTitulo}{azul}
\colorlet{azulObjetivos}{azul}
\colorlet{ifscGreen}{verdeVivo}
\colorlet{vermelhoCodigo}{vermelho}
\colorlet{cinzaCodigo}{fundo}
\colorlet{cinzaFundo}{fundo}
\colorlet{cinzaTabela}{fundo}
\colorlet{cinzaLinha}{fundo}
\colorlet{cinzaCabecalho}{barraTitulo}
\colorlet{pretoTabela}{borda}
\colorlet{borderGray}{borda}
\colorlet{bordaCodigo}{borda}
\colorlet{jsKeyword}{roxo}
\colorlet{jsString}{verde}
\colorlet{jsComment}{suave}
\colorlet{htmlTag}{azulMedio}
\colorlet{htmlString}{verde}
\colorlet{pyKeyword}{laranja}
\colorlet{pyBuiltin}{ciano}

\color{tinta}

% =====================================================
% LISTINGS — linguagens e estilos
% Regra de cópia: nenhuma linha de código pode passar de 103 caracteres
% (é o que cabe numa linha da janela de código; acima disso o PDF quebra a linha).
% =====================================================
\newcommand{\fonteCodigo}{\fontsize{8.5pt}{11.2pt}\selectfont\ttfamily}

\lstdefinelanguage{JavaScriptES6}{
  keywords={typeof, new, true, false, catch, function, return, null, switch, var, if, in, of, while, do, else, case, break, const, let, export, import, from, default, async, await, try, finally, throw, for, class, this},
  sensitive=true,
  comment=[l]{//},
  morecomment=[s]{/*}{*/},
  morestring=[b]',
  morestring=[b]",
  morestring=[b]`
}
\lstdefinelanguage{PythonCustom}{
  keywords={def, return, import, from, as, if, elif, else, while, for, in, try, except, finally, with, class, None, True, False, and, or, not, pass},
  ndkeywords={self, app, request, jsonify, CORS, sqlite3},
  sensitive=true,
  comment=[l]{\#},
  morestring=[b]',
  morestring=[b]",
  morestring=[s]{'''}{'''},
  morestring=[s]{"""}{"""}
}
\lstdefinelanguage{HTML5}{language=html}

\lstdefinestyle{base}{
    basicstyle=\fonteCodigo\color{tinta},
    keywordstyle=\color{roxo}\bfseries,
    ndkeywordstyle=\color{ciano}\bfseries,
    commentstyle=\color{suave}\itshape,
    stringstyle=\color{verde},
    identifierstyle=\color{tinta},
    breaklines=true, columns=fullflexible, keepspaces=true,
    showstringspaces=false, upquote=true, tabsize=2,
    frame=none, numbers=none, aboveskip=0pt, belowskip=0pt,
    extendedchars=true,
    literate=
      {á}{{\'a}}1 {é}{{\'e}}1 {í}{{\'i}}1 {ó}{{\'o}}1 {ú}{{\'u}}1
      {Á}{{\'A}}1 {É}{{\'E}}1 {Í}{{\'I}}1 {Ó}{{\'O}}1 {Ú}{{\'U}}1
      {à}{{\`a}}1 {è}{{\`e}}1 {ì}{{\`i}}1 {ò}{{\`o}}1 {ù}{{\`u}}1
      {ã}{{\~a}}1 {õ}{{\~o}}1 {Ã}{{\~A}}1 {Õ}{{\~O}}1
      {â}{{\^a}}1 {ê}{{\^e}}1 {î}{{\^i}}1 {ô}{{\^o}}1 {û}{{\^u}}1
      {Â}{{\^A}}1 {Ê}{{\^E}}1 {Ô}{{\^O}}1
      {ç}{{\c c}}1 {Ç}{{\c C}}1
      {–}{{--}}1 {—}{{---}}1 {±}{{$\pm$}}1 {≈}{{$\approx$}}1 {÷}{{$\div$}}1
}
\lstdefinestyle{js}{style=base, language=JavaScriptES6}
\lstdefinestyle{py}{style=base, language=PythonCustom, keywordstyle=\color{laranja}\bfseries, tabsize=4}
\lstdefinestyle{html}{style=base, language=HTML5, tagstyle=\color{azulMedio}\bfseries,
    keywordstyle=\color{azulMedio}, stringstyle=\color{verde}}
\lstdefinestyle{txt}{style=base, language={}}
\lstdefinestyle{saida}{style=base, language={}, basicstyle=\fonteCodigo\color{white!88!black},
    identifierstyle=\color{white!88!black}}
\lstset{style=base}

% =====================================================
% COMPONENTES VISUAIS
% =====================================================
% --- Janela de código (estilo editor) ---------------
% \begin{codigo}[js|html|py|txt]{nome-do-arquivo} ... \end{codigo}
\newcommand{\pontosJanela}{%
  {\color{vermelho!70}\tiny\faCircle}\,{\color{ambar!80}\tiny\faCircle}\,{\color{verdeVivo!70}\tiny\faCircle}\hspace{0.8em}}
\tcbset{janela/.style={
    enhanced, listing only, breakable,
    colback=fundo, colframe=borda, boxrule=0.6pt, arc=6pt,
    left=8pt, right=8pt, top=6pt, bottom=6pt,
    coltitle=suave, colbacktitle=barraTitulo,
    fonttitle=\ttfamily\footnotesize, toptitle=2pt, bottomtitle=2pt,
    titlerule=0.6pt, titlerule style=borda,
    before skip=6pt, after skip=10pt}}
% Etiqueta da linguagem, à direita da barra de título
\newcommand{\rotuloLinguagem}[2]{\hfill{\sffamily\scriptsize\bfseries\color{#1}#2}}
\newtcblisting{codigo}[2][js]{janela, title={\pontosJanela #2}, listing options={style=#1}}

% --- Janelas por linguagem (nomes usados nos roteiros) ----------------
% O argumento opcional é o nome do arquivo mostrado na barra da janela.
\newtcblisting{boxCodigoJS}[1][]{janela,
    title={\pontosJanela#1\rotuloLinguagem{ambar}{JavaScript}},
    listing options={style=js}}
\newtcblisting{boxCodigoHTML}[1][]{janela,
    title={\pontosJanela#1\rotuloLinguagem{vermelho}{HTML}},
    listing options={style=html}}
\newtcblisting{boxCodigoPY}[1][]{janela,
    title={\pontosJanela#1\rotuloLinguagem{azulMedio}{Python}},
    listing options={style=py}}
\newtcblisting{boxCodigoCSS}[1][]{janela,
    title={\pontosJanela#1\rotuloLinguagem{roxo}{CSS}},
    listing options={style=txt}}
% Croqui (desenho em texto da tela final)
\newtcblisting{boxCroqui}[1][]{janela,
    title={{\color{azulMedio}\faDraftingCompass}\hspace{0.5em}{\sffamily\bfseries #1}},
    listing options={style=txt}}

% Árvore de pastas do projeto
\newtcblisting{arvore}[1][Estrutura de pastas do projeto]{janela,
    title={{\color{ambar}\faFolderOpen}\hspace{0.5em}{\sffamily\bfseries #1}},
    listing options={style=txt}}

% --- Saída / terminal (fundo escuro) -------------------------------
\newtcblisting{saida}[1][Terminal]{
    enhanced, listing only, breakable,
    colback=terminal, colframe=terminal, boxrule=0pt, arc=6pt,
    left=10pt, right=10pt, top=4pt, bottom=6pt,
    title={\faTerminal\hspace{0.5em}#1}, coltitle=white!65!black,
    colbacktitle=terminal, fonttitle=\sffamily\footnotesize\bfseries,
    toptitle=3pt, bottomtitle=0pt, titlerule=0pt,
    before skip=6pt, after skip=10pt,
    listing options={style=saida}
}
% Versões antigas (conteúdo LaTeX com verbatim dentro)
\newtcolorbox{boxCodigoTerminal}{
    enhanced, breakable,
    colback=terminal, colframe=terminal, boxrule=0pt, arc=6pt, coltext=white!88!black,
    left=10pt, right=10pt, top=4pt, bottom=6pt,
    title={\faTerminal\hspace{0.5em}Terminal}, coltitle=white!65!black,
    colbacktitle=terminal, fonttitle=\sffamily\footnotesize\bfseries,
    toptitle=3pt, bottomtitle=0pt, titlerule=0pt,
    before upper={\singlespacing}
}
\newtcolorbox{boxCodigo}{
    enhanced, breakable,
    colback=fundo, colframe=borda, boxrule=0.6pt, arc=6pt,
    left=10pt, right=10pt, top=8pt, bottom=8pt,
    before upper={\singlespacing}
}

% --- Exemplo certo/errado -----------------------------
\newtcblisting{exemplo}[3][js]{
    enhanced, listing only,
    colback=#2!4, colframe=#2!45, boxrule=0.6pt, arc=6pt,
    left=8pt, right=8pt, top=5pt, bottom=5pt,
    title={#3}, colbacktitle=#2!12, coltitle=#2!75!black,
    fonttitle=\sffamily\small\bfseries, toptitle=2pt, bottomtitle=2pt,
    listing options={style=#1}
}

% --- Caixas de destaque com barra lateral -----------
\newtcolorbox{destaque}[3]{
    enhanced, breakable,
    colback=#1!6, frame hidden, arc=6pt,
    borderline west={3.5pt}{0pt}{#1},
    left=14pt, right=12pt, top=9pt, bottom=9pt,
    before skip=8pt, after skip=10pt,
    before upper={{\color{#1}\large #2}\hspace{0.55em}{\color{#1!75!black}\bfseries #3}\par\vspace{0.25em}}
}
\newenvironment{dica}[1]{\begin{destaque}{ambar}{\faLightbulb[regular]}{#1}}{\end{destaque}}
\newenvironment{atencao}[1]{\begin{destaque}{vermelho}{\faExclamationTriangle}{#1}}{\end{destaque}}
\newenvironment{conceito}[1]{\begin{destaque}{azulMedio}{\faBookOpen}{#1}}{\end{destaque}}
\newenvironment{sucesso}[1]{\begin{destaque}{verdeVivo}{\faCheckCircle[regular]}{#1}}{\end{destaque}}
\newenvironment{avancado}[1]{\begin{destaque}{roxo}{\faCode}{#1}}{\end{destaque}}
\newtcolorbox{repositorio}[1]{enhanced, colback=azul!5, frame hidden, arc=6pt,
    borderline west={3.5pt}{0pt}{azul}, left=14pt, right=12pt, top=9pt, bottom=9pt,
    before skip=8pt, after skip=10pt, fontupper=\small,
    before upper={{\color{azul}\large\faGithub}\hspace{0.55em}{\color{azul!75!black}\bfseries #1}\par\vspace{0.25em}}}
\newenvironment{checkpoint}[1]{\begin{destaque}{verdeVivo}{\faClipboardCheck}{#1}}{\end{destaque}}

% Nomes antigos
\newenvironment{boxDica}[1][Dica]{\begin{dica}{#1}}{\end{dica}}
\newenvironment{boxTeste}[1][Checkpoint de validação]{\begin{checkpoint}{#1}}{\end{checkpoint}}

% --- Caixa com título "aba" (objetivos, atividade) ---
\newtcolorbox{painel}[3][]{
    enhanced, breakable,
    colback=white, colframe=borda, boxrule=0.6pt, arc=8pt,
    left=14pt, right=14pt, top=14pt, bottom=10pt,
    title={#2\hspace{0.5em}#3}, coltitle=white, colbacktitle=azul,
    fonttitle=\bfseries, attach boxed title to top left={xshift=12pt, yshift=-8pt},
    boxed title style={arc=4pt, boxrule=0pt},
    before skip=12pt, after skip=10pt, #1
}
\newenvironment{boxObjetivos}[1][Ao final desta prática você será capaz de\ldots]
  {\begin{painel}{\faBullseye}{#1}}{\end{painel}}
\newenvironment{boxAtividade}[1][Atividade]
  {\begin{painel}[colbacktitle=azulMedio]{\faLaptopCode}{#1}}{\end{painel}}
\newtcolorbox{boxTabela}{
    enhanced, colback=white, colframe=borda, boxrule=0.6pt, arc=6pt,
    left=6pt, right=6pt, top=6pt, bottom=6pt}

% --- Checklist de entrega (itens com \caixaMarcar) ----------------
\newtcolorbox{checklist}[1][Antes de entregar, confira se\ldots]{enhanced, breakable,
    colback=verdeVivo!5, colframe=verdeVivo!40, boxrule=0.6pt, arc=8pt,
    left=14pt, right=14pt, top=10pt, bottom=10pt,
    title={\faClipboardCheck\hspace{0.5em}#1}, colbacktitle=verdeVivo, coltitle=white, fonttitle=\bfseries,
    before skip=10pt, after skip=10pt}

% --- Navegador (prévia da tela) -----------------------
\newtcolorbox{mockupNavegador}[1][]{
    enhanced, breakable,
    colback=white, colframe=borda, boxrule=0.8pt, arc=7pt,
    left=8pt, right=8pt, top=8pt, bottom=8pt,
    title={\pontosJanela\tcbox[on line, colback=white, colframe=borda, boxrule=0.5pt, arc=6pt,
        left=6pt, right=6pt, top=0.5pt, bottom=0.5pt, boxsep=0.5pt]{\ttfamily\scriptsize\color{suave}\faLock\ 127.0.0.1:5500/index.html}%
        \hfill{\sffamily\footnotesize\bfseries\color{suave}#1}},
    colbacktitle=barraTitulo, coltitle=tinta, titlerule=0.6pt, titlerule style=borda,
    toptitle=2pt, bottomtitle=2pt
}

% --- Cartões pequenos ---------------------------------
\tcbset{cartao/.style={enhanced, colback=fundo, colframe=borda, boxrule=0.6pt, arc=8pt,
    left=10pt, right=10pt, top=8pt, bottom=8pt, fontupper=\small\raggedright}}

% --- Elementos em linha --------------------------------
\newcommand{\tecla}[1]{\tcbox[on line, enhanced, colback=white, colframe=suave!50, boxrule=0.5pt,
    arc=2.5pt, left=3pt, right=3pt, top=0.5pt, bottom=0.5pt, boxsep=0.5pt,
    borderline south={1.2pt}{0pt}{suave!40}]{\ttfamily\footnotesize #1}}
\newcommand{\botao}[1]{\tcbox[on line, enhanced, colback=azulMedio!8, colframe=azulMedio!35, boxrule=0.5pt,
    arc=3pt, left=3pt, right=3pt, top=0.5pt, bottom=0.5pt, boxsep=0.5pt]{\sffamily\footnotesize\color{azul}#1}}
\newcommand{\seta}{\,{\color{suave}\scriptsize\faChevronRight}\,}
\newcommand{\chip}[2]{\tcbox[on line, enhanced, colback=white, opacityback=0.14, colframe=white,
    opacityframe=0.45, boxrule=0.5pt, arc=8pt, left=5pt, right=5pt, top=1.5pt, bottom=1.5pt,
    boxsep=0pt, coltext=white]{\footnotesize #1\hspace{0.4em}#2}}
\newcommand{\circulo}[2][azul]{\tikz[baseline=(c.base)]\node[circle, fill=#1, text=white,
    inner sep=0pt, minimum size=1.35em, font=\footnotesize\bfseries](c){#2};}
\newcommand{\marca}[1]{\tikz[baseline=(c.base)]\node[circle, fill=vermelho, text=white,
    inner sep=0pt, minimum size=1.2em, font=\scriptsize\bfseries](c){#1};}
\newcommand{\caixaMarcar}{{\color{suave}\faSquare[regular]}}

% --- Mantém um bloco de código curto inteiro na mesma página --------
% \manterJunto{n}: n = número de linhas do código que vem a seguir
\newcommand{\manterJunto}[1]{\needspace{\dimexpr 11.2pt*#1+5\baselineskip\relax}}

% --- Passo numerado do passo a passo ------------------
% \passo{1.2}{Criar a página HTML}  — o nome do arquivo vai na barra da janela de código
\newcommand{\passo}[2]{%
  \par\vspace{0.5em}\noindent
  \tcbox[on line, colback=azul, colframe=azul, arc=4pt, boxrule=0pt,
      left=5pt, right=5pt, top=1.5pt, bottom=1.5pt, boxsep=0pt]{\sffamily\footnotesize\bfseries\color{white}Passo #1}%
  \hspace{0.55em}{\bfseries\color{azul}#2}%
  \par\nopagebreak\vspace{-0.2em}}

% =====================================================
% ABERTURA PADRÃO DO ROTEIRO
% =====================================================
% \capaRoteiro{ícone-marca-d'água}{Título}{Subtítulo}{chips}
\newcommand{\capaRoteiro}[4]{%
\begin{tcolorbox}[enhanced, frame hidden, arc=12pt,
    interior style={left color=azul, right color=verde},
    left=20pt, right=20pt, top=18pt, bottom=16pt, colupper=white,
    overlay={\node[anchor=east, opacity=0.13, font=\fontsize{90}{90}\selectfont, text=white] at ([xshift=-14pt]frame.east) {#1};}]
{\footnotesize\bfseries\color{white!80!azul}ROTEIRO PRÁTICO \numeroRoteiro\hspace{0.6em}\textbullet\hspace{0.6em}PROGRAMAÇÃO PARA A INTERNET 2}\par\vspace{2pt}
{\fontsize{24}{29}\selectfont\bfseries #2}\par\vspace{2pt}
{\large #3}\par\vspace{10pt}
#4
\end{tcolorbox}}

% \cartoesUC{Produto/entrega deste roteiro}
\newcommand{\cartoesUC}[1]{%
\begin{tcbraster}[raster columns=3, raster equal height, raster column skip=8pt]
\begin{tcolorbox}[cartao]
{\color{azulMedio}\faUniversity}\hspace{0.4em}\textbf{Unidade Curricular}\par\vspace{2pt}
Programação para a Internet 2\\ {\color{suave}CST em Sistemas para a Internet}
\end{tcolorbox}
\begin{tcolorbox}[cartao]
{\color{azulMedio}\faChalkboardTeacher}\hspace{0.4em}\textbf{Docente}\par\vspace{2pt}
Prof. André Moraes\\ {\color{suave}IFSC Câmpus Garopaba}
\end{tcolorbox}
\begin{tcolorbox}[cartao]
{\color{azulMedio}\faFlag}\hspace{0.4em}\textbf{Produto deste roteiro}\par\vspace{2pt}
#1
\end{tcolorbox}
\end{tcbraster}}

% \etapaMapa{cor}{n}{Título}{descrição}  — usar dentro de \begin{mapaAula} ... \end{mapaAula}
\newenvironment{mapaAula}{%
  \needspace{9\baselineskip}{\large\bfseries\color{azul}\faMap[regular]\hspace{0.5em}Mapa da aula}\par\vspace{0.2em}
  \begin{tcbraster}[raster columns=2, raster equal height, raster column skip=8pt, raster row skip=6pt]}
  {\end{tcbraster}}
\newcommand{\etapaMapa}[4]{%
\begin{tcolorbox}[cartao, colback=#1!5, colframe=#1!25]
\circulo[#1]{#2}\hspace{0.4em}\textbf{#3}\par\vspace{2pt}
{\color{suave}\footnotesize #4}
\end{tcolorbox}}

% \trilhaRoteiros{n} — trilha da disciplina com o roteiro atual em destaque
\newcommand{\trilhaRoteiros}[1]{%
\begin{center}
\begin{tikzpicture}[node distance=0.32cm,
    etapa/.style={rectangle, draw=borda, fill=fundo, line width=0.6pt, text width=2.55cm, align=center,
        rounded corners=7pt, minimum height=2.0cm, font=\footnotesize},
    atual/.style={draw=verdeVivo, fill=verdeVivo!7, line width=1.1pt},
    linha/.style={draw=suave!70, -{Latex[length=2.2mm]}, line width=0.9pt},
    r1/.style={}, r2/.style={}, r3/.style={}, r4/.style={}, r5/.style={}, r#1/.style={atual}]
    \node[etapa, r1] (r1) {{\color{azulMedio}\Large\faColumns}\\[3pt]\textbf{Roteiro 01}\\{\color{suave}Bootstrap 5 e layout}};
    \node[etapa, right=of r1, r2] (r2) {{\color{azulMedio}\Large\faMousePointer}\\[3pt]\textbf{Roteiro 02}\\{\color{suave}jQuery e DOM}};
    \node[etapa, right=of r2, r3] (r3) {{\color{azulMedio}\Large\faCalculator}\\[3pt]\textbf{Roteiro 03}\\{\color{suave}Calculadora e LocalStorage}};
    \node[etapa, right=of r3, r4] (r4) {{\color{azulMedio}\Large\faLayerGroup}\\[3pt]\textbf{Roteiro 04}\\{\color{suave}ES6 Modules em camadas}};
    \node[etapa, right=of r4, r5] (r5) {{\color{azulMedio}\Large\faServer}\\[3pt]\textbf{Roteiro 05}\\{\color{suave}Fetch API e Flask}};
    \draw[linha] (r1) -- (r2);
    \draw[linha] (r2) -- (r3);
    \draw[linha] (r3) -- (r4);
    \draw[linha] (r4) -- (r5);
\end{tikzpicture}
\end{center}}

% =====================================================
% HIPERLINKS
% =====================================================
\hypersetup{colorlinks=true, linkcolor=azulMedio, urlcolor=verde, citecolor=azul}

% =====================================================
% CABEÇALHO E RODAPÉ
% =====================================================
\pagestyle{fancy}
\fancyhf{}
\fancyhead[L]{%
    \IfFileExists{../figures/garopaba_horizontal_marca2015_PNG.png}{%
        \raisebox{-0.25\height}{\includegraphics[height=0.95cm]{../figures/garopaba_horizontal_marca2015_PNG.png}}%
    }{\textbf{\color{verde} IFSC — Câmpus Garopaba}}}
\fancyhead[R]{\footnotesize\color{suave}\textbf{\color{azul}Roteiro Prático \numeroRoteiro}\hspace{0.5em}\textbullet\hspace{0.5em}\tituloCurto}
\fancyfoot[L]{\footnotesize\color{suave}Programação para a Internet 2\hspace{0.4em}\textbullet\hspace{0.4em}CST Sistemas para a Internet}
\fancyfoot[R]{\tcbox[on line, colback=azul!8, colframe=azul!8, arc=6pt, boxrule=0pt, left=5pt, right=5pt, top=1pt, bottom=1pt, boxsep=0pt]{\footnotesize\color{azul}\bfseries\thepage\,/\,\pageref*{LastPage}}}
\renewcommand{\headrule}{{\color{borda}\hrule height 0.6pt}}
\renewcommand{\footrulewidth}{0pt}
\setlength{\headheight}{36pt}

% =====================================================
% TÍTULOS DE SEÇÃO
% =====================================================
\titleformat{\section}
  {\normalfont\Large\bfseries\color{azul}}
  {\tikz[baseline=(n.base)]\node[fill=azul, rounded corners=5pt, text=white, inner sep=5pt,
      minimum width=1.9em, font=\large\bfseries](n){\thesection};}
  {0.7em}{}
\titlespacing*{\section}{0pt}{1.6em}{0.7em}

\titleformat{\subsection}
  {\normalfont\large\bfseries\color{azul!85!black}}
  {{\color{verdeVivo}\thesubsection}}
  {0.6em}{}
\titlespacing*{\subsection}{0pt}{1.2em}{0.4em}

\titleformat{\subsubsection}
  {\normalfont\normalsize\bfseries\color{azul!85!black}}
  {{\color{verdeVivo}\thesubsubsection}}
  {0.6em}{}
\titlespacing*{\subsubsection}{0pt}{1.0em}{0.3em}

% Títulos nunca ficam sozinhos no pé da página
\pretocmd{\section}{\needspace{7\baselineskip}}{}{}
\pretocmd{\subsection}{\needspace{5\baselineskip}}{}{}

\setlist[itemize]{leftmargin=1.4em, itemsep=0.25em, topsep=0.3em, label={\color{verdeVivo}\small\faAngleRight}}
\setlist[enumerate]{leftmargin=2.2em, itemsep=0.45em, topsep=0.3em}

% Tabelas: sem linhas verticais, cabeçalho colorido, zebra suave
\newcommand{\cabecalho}{\rowcolor{azul}}
\newcommand{\cab}[1]{\leavevmode\color{white}\bfseries #1}
\arrayrulecolor{borda}
\newcolumntype{Y}{>{\raggedright\arraybackslash}X}
\newcolumntype{P}[1]{>{\raggedright\arraybackslash}p{#1}}
% Primeira coluna em destaque (rótulo da linha)
\newcolumntype{R}{>{\raggedright\arraybackslash\bfseries\leavevmode\color{azul}}l}
\newcolumntype{Q}[1]{>{\raggedright\arraybackslash\bfseries\leavevmode\color{azul}}p{#1}}

% Legendas de figuras e tabelas discretas
\usepackage[font={small,color=suave}, labelfont={bf,color=azul}, skip=6pt]{caption}

% \begin{document} ... (conteúdo do roteiro)

%   \begin{document}
%   \input{content.tex}
%   \end{document}
% Compilar com: tectonic main.tex
% =====================================================================
\documentclass[11pt,a4paper]{article}

\providecommand{\numeroRoteiro}{0X}
\providecommand{\tituloCurto}{Título curto do roteiro}

% =====================================================
% PACOTES
% =====================================================
\usepackage[utf8]{inputenc}
\usepackage[T1]{fontenc}
\usepackage[brazil]{babel}
\shorthandoff{"}
\AtBeginDocument{\shorthandoff{"}}
\usepackage[default]{sourcesanspro}          % texto: sem serifa, leitura confortável
\usepackage[scale=0.86]{sourcecodepro}       % código: monoespaçada moderna
% =====================================================
% ÍCONES (fontawesome5 com compatibilidade Tectonic no macOS)
% Evita crash de alocação na rotina nativa \XeTeXglyphname do Tectonic
% =====================================================
\usepackage{expl3}
\ExplSyntaxOn
\cs_set:Npn \tex_XeTeXglyphname:D #1#2 { }
\ExplSyntaxOff
\newcommand{\faDice}{}
\usepackage{fontawesome5}

\NewDocumentCommand\definirFaIcone{m m}{%
  \expandafter\NewDocumentCommand\csname fa#1\endcsname{O{}}{%
    \if\relax\detokenize{##1}\relax
      \faIcon{#2}%
    \else
      \faIcon[##1]{#2}%
    \fi
  }%
}
\definirFaIcone{AngleRight}{angle-right}
\definirFaIcone{BookOpen}{book-open}
\definirFaIcone{Bullseye}{bullseye}
\definirFaIcone{Calculator}{calculator}
\definirFaIcone{ChalkboardTeacher}{chalkboard-teacher}
\definirFaIcone{CheckCircle}{check-circle}
\definirFaIcone{ChevronRight}{chevron-right}
\definirFaIcone{Circle}{circle}
\definirFaIcone{ClipboardCheck}{clipboard-check}
\definirFaIcone{Code}{code}
\definirFaIcone{Columns}{columns}
\definirFaIcone{Database}{database}
\definirFaIcone{DraftingCompass}{drafting-compass}
\definirFaIcone{ExclamationTriangle}{exclamation-triangle}
\definirFaIcone{ExternalLinkAlt}{external-link-alt}
\definirFaIcone{Flag}{flag}
\definirFaIcone{FolderOpen}{folder-open}
\definirFaIcone{Github}{github}
\definirFaIcone{LaptopCode}{laptop-code}
\definirFaIcone{LayerGroup}{layer-group}
\definirFaIcone{Lightbulb}{lightbulb}
\definirFaIcone{Link}{link}
\definirFaIcone{Lock}{lock}
\definirFaIcone{Map}{map}
\definirFaIcone{MapMarkerAlt}{map-marker-alt}
\definirFaIcone{MousePointer}{mouse-pointer}
\definirFaIcone{Server}{server}
\definirFaIcone{Square}{square}
\definirFaIcone{Terminal}{terminal}
\definirFaIcone{Tools}{tools}
\definirFaIcone{University}{university}
\definirFaIcone{Check}{check}
\definirFaIcone{Times}{times}
\definirFaIcone{InfoCircle}{info-circle}
\definirFaIcone{QuestionCircle}{question-circle}
\definirFaIcone{Play}{play}
\definirFaIcone{FileCode}{file-code}
\definirFaIcone{FileAlt}{file-alt}
\definirFaIcone{Folder}{folder}
\definirFaIcone{Cogs}{cogs}
\definirFaIcone{Cog}{cog}
\definirFaIcone{Eye}{eye}
\definirFaIcone{Desktop}{desktop}
\definirFaIcone{Globe}{globe}
\definirFaIcone{Search}{search}
\definirFaIcone{Download}{download}
\definirFaIcone{Upload}{upload}
\definirFaIcone{Sync}{sync}
\definirFaIcone{ShareAlt}{share-alt}
\definirFaIcone{Edit}{edit}
\definirFaIcone{Trash}{trash}
\definirFaIcone{Plus}{plus}
\definirFaIcone{Minus}{minus}
\usepackage{geometry}
\usepackage{graphicx}
\graphicspath{{../figures/}{figures/}}
\usepackage{fancyhdr}
\usepackage{titlesec}
\usepackage[table]{xcolor}
\usepackage{hyperref}
\usepackage{xurl}                            % URLs longas quebram em qualquer ponto
\usepackage{enumitem}
\usepackage[most]{tcolorbox}
\usepackage{float}
\usepackage{tikz}
\usetikzlibrary{arrows.meta, positioning, shapes.geometric, calc, babel}
\usepackage{tabularx}
\usepackage{booktabs}
\usepackage{listings}
\usepackage{upquote}
\usepackage{lastpage}
\usepackage{qrcode}
\usepackage{multicol}
\usepackage{setspace}
\usepackage{needspace}
\usepackage{adjustbox}

\geometry{top=2.6cm, bottom=2.6cm, left=2cm, right=2cm, headsep=0.8cm}
\linespread{1.18}
\setlength{\parskip}{0.55em}
\setlength{\parindent}{0pt}
\setlength{\emergencystretch}{6em}   % evita linhas estouradas com \texttt longos

% =====================================================
% PALETA (cores institucionais IFSC, em tons suaves)
% =====================================================
\definecolor{azul}{RGB}{30,60,90}
\definecolor{azulMedio}{RGB}{37,99,235}
\definecolor{verde}{RGB}{22,101,52}
\definecolor{verdeVivo}{RGB}{22,163,74}
\definecolor{ambar}{RGB}{217,119,6}
\definecolor{vermelho}{RGB}{220,38,38}
\definecolor{roxo}{RGB}{124,58,237}
\definecolor{laranja}{RGB}{194,65,12}
\definecolor{ciano}{RGB}{14,116,144}
\definecolor{tinta}{RGB}{30,41,59}
\definecolor{suave}{RGB}{100,116,139}
\definecolor{borda}{RGB}{226,232,240}
\definecolor{fundo}{RGB}{248,250,252}
\definecolor{barraTitulo}{RGB}{241,245,249}
\definecolor{terminal}{RGB}{15,23,42}

% Nomes antigos dos roteiros de ProgWeb 2 (mantidos para compatibilidade)
\colorlet{azulTitulo}{azul}
\colorlet{azulObjetivos}{azul}
\colorlet{ifscGreen}{verdeVivo}
\colorlet{vermelhoCodigo}{vermelho}
\colorlet{cinzaCodigo}{fundo}
\colorlet{cinzaFundo}{fundo}
\colorlet{cinzaTabela}{fundo}
\colorlet{cinzaLinha}{fundo}
\colorlet{cinzaCabecalho}{barraTitulo}
\colorlet{pretoTabela}{borda}
\colorlet{borderGray}{borda}
\colorlet{bordaCodigo}{borda}
\colorlet{jsKeyword}{roxo}
\colorlet{jsString}{verde}
\colorlet{jsComment}{suave}
\colorlet{htmlTag}{azulMedio}
\colorlet{htmlString}{verde}
\colorlet{pyKeyword}{laranja}
\colorlet{pyBuiltin}{ciano}

\color{tinta}

% =====================================================
% LISTINGS — linguagens e estilos
% Regra de cópia: nenhuma linha de código pode passar de 103 caracteres
% (é o que cabe numa linha da janela de código; acima disso o PDF quebra a linha).
% =====================================================
\newcommand{\fonteCodigo}{\fontsize{8.5pt}{11.2pt}\selectfont\ttfamily}

\lstdefinelanguage{JavaScriptES6}{
  keywords={typeof, new, true, false, catch, function, return, null, switch, var, if, in, of, while, do, else, case, break, const, let, export, import, from, default, async, await, try, finally, throw, for, class, this},
  sensitive=true,
  comment=[l]{//},
  morecomment=[s]{/*}{*/},
  morestring=[b]',
  morestring=[b]",
  morestring=[b]`
}
\lstdefinelanguage{PythonCustom}{
  keywords={def, return, import, from, as, if, elif, else, while, for, in, try, except, finally, with, class, None, True, False, and, or, not, pass},
  ndkeywords={self, app, request, jsonify, CORS, sqlite3},
  sensitive=true,
  comment=[l]{\#},
  morestring=[b]',
  morestring=[b]",
  morestring=[s]{'''}{'''},
  morestring=[s]{"""}{"""}
}
\lstdefinelanguage{HTML5}{language=html}

\lstdefinestyle{base}{
    basicstyle=\fonteCodigo\color{tinta},
    keywordstyle=\color{roxo}\bfseries,
    ndkeywordstyle=\color{ciano}\bfseries,
    commentstyle=\color{suave}\itshape,
    stringstyle=\color{verde},
    identifierstyle=\color{tinta},
    breaklines=true, columns=fullflexible, keepspaces=true,
    showstringspaces=false, upquote=true, tabsize=2,
    frame=none, numbers=none, aboveskip=0pt, belowskip=0pt,
    extendedchars=true,
    literate=
      {á}{{\'a}}1 {é}{{\'e}}1 {í}{{\'i}}1 {ó}{{\'o}}1 {ú}{{\'u}}1
      {Á}{{\'A}}1 {É}{{\'E}}1 {Í}{{\'I}}1 {Ó}{{\'O}}1 {Ú}{{\'U}}1
      {à}{{\`a}}1 {è}{{\`e}}1 {ì}{{\`i}}1 {ò}{{\`o}}1 {ù}{{\`u}}1
      {ã}{{\~a}}1 {õ}{{\~o}}1 {Ã}{{\~A}}1 {Õ}{{\~O}}1
      {â}{{\^a}}1 {ê}{{\^e}}1 {î}{{\^i}}1 {ô}{{\^o}}1 {û}{{\^u}}1
      {Â}{{\^A}}1 {Ê}{{\^E}}1 {Ô}{{\^O}}1
      {ç}{{\c c}}1 {Ç}{{\c C}}1
      {–}{{--}}1 {—}{{---}}1 {±}{{$\pm$}}1 {≈}{{$\approx$}}1 {÷}{{$\div$}}1
}
\lstdefinestyle{js}{style=base, language=JavaScriptES6}
\lstdefinestyle{py}{style=base, language=PythonCustom, keywordstyle=\color{laranja}\bfseries, tabsize=4}
\lstdefinestyle{html}{style=base, language=HTML5, tagstyle=\color{azulMedio}\bfseries,
    keywordstyle=\color{azulMedio}, stringstyle=\color{verde}}
\lstdefinestyle{txt}{style=base, language={}}
\lstdefinestyle{saida}{style=base, language={}, basicstyle=\fonteCodigo\color{white!88!black},
    identifierstyle=\color{white!88!black}}
\lstset{style=base}

% =====================================================
% COMPONENTES VISUAIS
% =====================================================
% --- Janela de código (estilo editor) ---------------
% \begin{codigo}[js|html|py|txt]{nome-do-arquivo} ... \end{codigo}
\newcommand{\pontosJanela}{%
  {\color{vermelho!70}\tiny\faCircle}\,{\color{ambar!80}\tiny\faCircle}\,{\color{verdeVivo!70}\tiny\faCircle}\hspace{0.8em}}
\tcbset{janela/.style={
    enhanced, listing only, breakable,
    colback=fundo, colframe=borda, boxrule=0.6pt, arc=6pt,
    left=8pt, right=8pt, top=6pt, bottom=6pt,
    coltitle=suave, colbacktitle=barraTitulo,
    fonttitle=\ttfamily\footnotesize, toptitle=2pt, bottomtitle=2pt,
    titlerule=0.6pt, titlerule style=borda,
    before skip=6pt, after skip=10pt}}
% Etiqueta da linguagem, à direita da barra de título
\newcommand{\rotuloLinguagem}[2]{\hfill{\sffamily\scriptsize\bfseries\color{#1}#2}}
\newtcblisting{codigo}[2][js]{janela, title={\pontosJanela #2}, listing options={style=#1}}

% --- Janelas por linguagem (nomes usados nos roteiros) ----------------
% O argumento opcional é o nome do arquivo mostrado na barra da janela.
\newtcblisting{boxCodigoJS}[1][]{janela,
    title={\pontosJanela#1\rotuloLinguagem{ambar}{JavaScript}},
    listing options={style=js}}
\newtcblisting{boxCodigoHTML}[1][]{janela,
    title={\pontosJanela#1\rotuloLinguagem{vermelho}{HTML}},
    listing options={style=html}}
\newtcblisting{boxCodigoPY}[1][]{janela,
    title={\pontosJanela#1\rotuloLinguagem{azulMedio}{Python}},
    listing options={style=py}}
\newtcblisting{boxCodigoCSS}[1][]{janela,
    title={\pontosJanela#1\rotuloLinguagem{roxo}{CSS}},
    listing options={style=txt}}
% Croqui (desenho em texto da tela final)
\newtcblisting{boxCroqui}[1][]{janela,
    title={{\color{azulMedio}\faDraftingCompass}\hspace{0.5em}{\sffamily\bfseries #1}},
    listing options={style=txt}}

% Árvore de pastas do projeto
\newtcblisting{arvore}[1][Estrutura de pastas do projeto]{janela,
    title={{\color{ambar}\faFolderOpen}\hspace{0.5em}{\sffamily\bfseries #1}},
    listing options={style=txt}}

% --- Saída / terminal (fundo escuro) -------------------------------
\newtcblisting{saida}[1][Terminal]{
    enhanced, listing only, breakable,
    colback=terminal, colframe=terminal, boxrule=0pt, arc=6pt,
    left=10pt, right=10pt, top=4pt, bottom=6pt,
    title={\faTerminal\hspace{0.5em}#1}, coltitle=white!65!black,
    colbacktitle=terminal, fonttitle=\sffamily\footnotesize\bfseries,
    toptitle=3pt, bottomtitle=0pt, titlerule=0pt,
    before skip=6pt, after skip=10pt,
    listing options={style=saida}
}
% Versões antigas (conteúdo LaTeX com verbatim dentro)
\newtcolorbox{boxCodigoTerminal}{
    enhanced, breakable,
    colback=terminal, colframe=terminal, boxrule=0pt, arc=6pt, coltext=white!88!black,
    left=10pt, right=10pt, top=4pt, bottom=6pt,
    title={\faTerminal\hspace{0.5em}Terminal}, coltitle=white!65!black,
    colbacktitle=terminal, fonttitle=\sffamily\footnotesize\bfseries,
    toptitle=3pt, bottomtitle=0pt, titlerule=0pt,
    before upper={\singlespacing}
}
\newtcolorbox{boxCodigo}{
    enhanced, breakable,
    colback=fundo, colframe=borda, boxrule=0.6pt, arc=6pt,
    left=10pt, right=10pt, top=8pt, bottom=8pt,
    before upper={\singlespacing}
}

% --- Exemplo certo/errado -----------------------------
\newtcblisting{exemplo}[3][js]{
    enhanced, listing only,
    colback=#2!4, colframe=#2!45, boxrule=0.6pt, arc=6pt,
    left=8pt, right=8pt, top=5pt, bottom=5pt,
    title={#3}, colbacktitle=#2!12, coltitle=#2!75!black,
    fonttitle=\sffamily\small\bfseries, toptitle=2pt, bottomtitle=2pt,
    listing options={style=#1}
}

% --- Caixas de destaque com barra lateral -----------
\newtcolorbox{destaque}[3]{
    enhanced, breakable,
    colback=#1!6, frame hidden, arc=6pt,
    borderline west={3.5pt}{0pt}{#1},
    left=14pt, right=12pt, top=9pt, bottom=9pt,
    before skip=8pt, after skip=10pt,
    before upper={{\color{#1}\large #2}\hspace{0.55em}{\color{#1!75!black}\bfseries #3}\par\vspace{0.25em}}
}
\newenvironment{dica}[1]{\begin{destaque}{ambar}{\faLightbulb[regular]}{#1}}{\end{destaque}}
\newenvironment{atencao}[1]{\begin{destaque}{vermelho}{\faExclamationTriangle}{#1}}{\end{destaque}}
\newenvironment{conceito}[1]{\begin{destaque}{azulMedio}{\faBookOpen}{#1}}{\end{destaque}}
\newenvironment{sucesso}[1]{\begin{destaque}{verdeVivo}{\faCheckCircle[regular]}{#1}}{\end{destaque}}
\newenvironment{avancado}[1]{\begin{destaque}{roxo}{\faCode}{#1}}{\end{destaque}}
\newtcolorbox{repositorio}[1]{enhanced, colback=azul!5, frame hidden, arc=6pt,
    borderline west={3.5pt}{0pt}{azul}, left=14pt, right=12pt, top=9pt, bottom=9pt,
    before skip=8pt, after skip=10pt, fontupper=\small,
    before upper={{\color{azul}\large\faGithub}\hspace{0.55em}{\color{azul!75!black}\bfseries #1}\par\vspace{0.25em}}}
\newenvironment{checkpoint}[1]{\begin{destaque}{verdeVivo}{\faClipboardCheck}{#1}}{\end{destaque}}

% Nomes antigos
\newenvironment{boxDica}[1][Dica]{\begin{dica}{#1}}{\end{dica}}
\newenvironment{boxTeste}[1][Checkpoint de validação]{\begin{checkpoint}{#1}}{\end{checkpoint}}

% --- Caixa com título "aba" (objetivos, atividade) ---
\newtcolorbox{painel}[3][]{
    enhanced, breakable,
    colback=white, colframe=borda, boxrule=0.6pt, arc=8pt,
    left=14pt, right=14pt, top=14pt, bottom=10pt,
    title={#2\hspace{0.5em}#3}, coltitle=white, colbacktitle=azul,
    fonttitle=\bfseries, attach boxed title to top left={xshift=12pt, yshift=-8pt},
    boxed title style={arc=4pt, boxrule=0pt},
    before skip=12pt, after skip=10pt, #1
}
\newenvironment{boxObjetivos}[1][Ao final desta prática você será capaz de\ldots]
  {\begin{painel}{\faBullseye}{#1}}{\end{painel}}
\newenvironment{boxAtividade}[1][Atividade]
  {\begin{painel}[colbacktitle=azulMedio]{\faLaptopCode}{#1}}{\end{painel}}
\newtcolorbox{boxTabela}{
    enhanced, colback=white, colframe=borda, boxrule=0.6pt, arc=6pt,
    left=6pt, right=6pt, top=6pt, bottom=6pt}

% --- Checklist de entrega (itens com \caixaMarcar) ----------------
\newtcolorbox{checklist}[1][Antes de entregar, confira se\ldots]{enhanced, breakable,
    colback=verdeVivo!5, colframe=verdeVivo!40, boxrule=0.6pt, arc=8pt,
    left=14pt, right=14pt, top=10pt, bottom=10pt,
    title={\faClipboardCheck\hspace{0.5em}#1}, colbacktitle=verdeVivo, coltitle=white, fonttitle=\bfseries,
    before skip=10pt, after skip=10pt}

% --- Navegador (prévia da tela) -----------------------
\newtcolorbox{mockupNavegador}[1][]{
    enhanced, breakable,
    colback=white, colframe=borda, boxrule=0.8pt, arc=7pt,
    left=8pt, right=8pt, top=8pt, bottom=8pt,
    title={\pontosJanela\tcbox[on line, colback=white, colframe=borda, boxrule=0.5pt, arc=6pt,
        left=6pt, right=6pt, top=0.5pt, bottom=0.5pt, boxsep=0.5pt]{\ttfamily\scriptsize\color{suave}\faLock\ 127.0.0.1:5500/index.html}%
        \hfill{\sffamily\footnotesize\bfseries\color{suave}#1}},
    colbacktitle=barraTitulo, coltitle=tinta, titlerule=0.6pt, titlerule style=borda,
    toptitle=2pt, bottomtitle=2pt
}

% --- Cartões pequenos ---------------------------------
\tcbset{cartao/.style={enhanced, colback=fundo, colframe=borda, boxrule=0.6pt, arc=8pt,
    left=10pt, right=10pt, top=8pt, bottom=8pt, fontupper=\small\raggedright}}

% --- Elementos em linha --------------------------------
\newcommand{\tecla}[1]{\tcbox[on line, enhanced, colback=white, colframe=suave!50, boxrule=0.5pt,
    arc=2.5pt, left=3pt, right=3pt, top=0.5pt, bottom=0.5pt, boxsep=0.5pt,
    borderline south={1.2pt}{0pt}{suave!40}]{\ttfamily\footnotesize #1}}
\newcommand{\botao}[1]{\tcbox[on line, enhanced, colback=azulMedio!8, colframe=azulMedio!35, boxrule=0.5pt,
    arc=3pt, left=3pt, right=3pt, top=0.5pt, bottom=0.5pt, boxsep=0.5pt]{\sffamily\footnotesize\color{azul}#1}}
\newcommand{\seta}{\,{\color{suave}\scriptsize\faChevronRight}\,}
\newcommand{\chip}[2]{\tcbox[on line, enhanced, colback=white, opacityback=0.14, colframe=white,
    opacityframe=0.45, boxrule=0.5pt, arc=8pt, left=5pt, right=5pt, top=1.5pt, bottom=1.5pt,
    boxsep=0pt, coltext=white]{\footnotesize #1\hspace{0.4em}#2}}
\newcommand{\circulo}[2][azul]{\tikz[baseline=(c.base)]\node[circle, fill=#1, text=white,
    inner sep=0pt, minimum size=1.35em, font=\footnotesize\bfseries](c){#2};}
\newcommand{\marca}[1]{\tikz[baseline=(c.base)]\node[circle, fill=vermelho, text=white,
    inner sep=0pt, minimum size=1.2em, font=\scriptsize\bfseries](c){#1};}
\newcommand{\caixaMarcar}{{\color{suave}\faSquare[regular]}}

% --- Mantém um bloco de código curto inteiro na mesma página --------
% \manterJunto{n}: n = número de linhas do código que vem a seguir
\newcommand{\manterJunto}[1]{\needspace{\dimexpr 11.2pt*#1+5\baselineskip\relax}}

% --- Passo numerado do passo a passo ------------------
% \passo{1.2}{Criar a página HTML}  — o nome do arquivo vai na barra da janela de código
\newcommand{\passo}[2]{%
  \par\vspace{0.5em}\noindent
  \tcbox[on line, colback=azul, colframe=azul, arc=4pt, boxrule=0pt,
      left=5pt, right=5pt, top=1.5pt, bottom=1.5pt, boxsep=0pt]{\sffamily\footnotesize\bfseries\color{white}Passo #1}%
  \hspace{0.55em}{\bfseries\color{azul}#2}%
  \par\nopagebreak\vspace{-0.2em}}

% =====================================================
% ABERTURA PADRÃO DO ROTEIRO
% =====================================================
% \capaRoteiro{ícone-marca-d'água}{Título}{Subtítulo}{chips}
\newcommand{\capaRoteiro}[4]{%
\begin{tcolorbox}[enhanced, frame hidden, arc=12pt,
    interior style={left color=azul, right color=verde},
    left=20pt, right=20pt, top=18pt, bottom=16pt, colupper=white,
    overlay={\node[anchor=east, opacity=0.13, font=\fontsize{90}{90}\selectfont, text=white] at ([xshift=-14pt]frame.east) {#1};}]
{\footnotesize\bfseries\color{white!80!azul}ROTEIRO PRÁTICO \numeroRoteiro\hspace{0.6em}\textbullet\hspace{0.6em}PROGRAMAÇÃO PARA A INTERNET 2}\par\vspace{2pt}
{\fontsize{24}{29}\selectfont\bfseries #2}\par\vspace{2pt}
{\large #3}\par\vspace{10pt}
#4
\end{tcolorbox}}

% \cartoesUC{Produto/entrega deste roteiro}
\newcommand{\cartoesUC}[1]{%
\begin{tcbraster}[raster columns=3, raster equal height, raster column skip=8pt]
\begin{tcolorbox}[cartao]
{\color{azulMedio}\faUniversity}\hspace{0.4em}\textbf{Unidade Curricular}\par\vspace{2pt}
Programação para a Internet 2\\ {\color{suave}CST em Sistemas para a Internet}
\end{tcolorbox}
\begin{tcolorbox}[cartao]
{\color{azulMedio}\faChalkboardTeacher}\hspace{0.4em}\textbf{Docente}\par\vspace{2pt}
Prof. André Moraes\\ {\color{suave}IFSC Câmpus Garopaba}
\end{tcolorbox}
\begin{tcolorbox}[cartao]
{\color{azulMedio}\faFlag}\hspace{0.4em}\textbf{Produto deste roteiro}\par\vspace{2pt}
#1
\end{tcolorbox}
\end{tcbraster}}

% \etapaMapa{cor}{n}{Título}{descrição}  — usar dentro de \begin{mapaAula} ... \end{mapaAula}
\newenvironment{mapaAula}{%
  \needspace{9\baselineskip}{\large\bfseries\color{azul}\faMap[regular]\hspace{0.5em}Mapa da aula}\par\vspace{0.2em}
  \begin{tcbraster}[raster columns=2, raster equal height, raster column skip=8pt, raster row skip=6pt]}
  {\end{tcbraster}}
\newcommand{\etapaMapa}[4]{%
\begin{tcolorbox}[cartao, colback=#1!5, colframe=#1!25]
\circulo[#1]{#2}\hspace{0.4em}\textbf{#3}\par\vspace{2pt}
{\color{suave}\footnotesize #4}
\end{tcolorbox}}

% \trilhaRoteiros{n} — trilha da disciplina com o roteiro atual em destaque
\newcommand{\trilhaRoteiros}[1]{%
\begin{center}
\begin{tikzpicture}[node distance=0.32cm,
    etapa/.style={rectangle, draw=borda, fill=fundo, line width=0.6pt, text width=2.55cm, align=center,
        rounded corners=7pt, minimum height=2.0cm, font=\footnotesize},
    atual/.style={draw=verdeVivo, fill=verdeVivo!7, line width=1.1pt},
    linha/.style={draw=suave!70, -{Latex[length=2.2mm]}, line width=0.9pt},
    r1/.style={}, r2/.style={}, r3/.style={}, r4/.style={}, r5/.style={}, r#1/.style={atual}]
    \node[etapa, r1] (r1) {{\color{azulMedio}\Large\faColumns}\\[3pt]\textbf{Roteiro 01}\\{\color{suave}Bootstrap 5 e layout}};
    \node[etapa, right=of r1, r2] (r2) {{\color{azulMedio}\Large\faMousePointer}\\[3pt]\textbf{Roteiro 02}\\{\color{suave}jQuery e DOM}};
    \node[etapa, right=of r2, r3] (r3) {{\color{azulMedio}\Large\faCalculator}\\[3pt]\textbf{Roteiro 03}\\{\color{suave}Calculadora e LocalStorage}};
    \node[etapa, right=of r3, r4] (r4) {{\color{azulMedio}\Large\faLayerGroup}\\[3pt]\textbf{Roteiro 04}\\{\color{suave}ES6 Modules em camadas}};
    \node[etapa, right=of r4, r5] (r5) {{\color{azulMedio}\Large\faServer}\\[3pt]\textbf{Roteiro 05}\\{\color{suave}Fetch API e Flask}};
    \draw[linha] (r1) -- (r2);
    \draw[linha] (r2) -- (r3);
    \draw[linha] (r3) -- (r4);
    \draw[linha] (r4) -- (r5);
\end{tikzpicture}
\end{center}}

% =====================================================
% HIPERLINKS
% =====================================================
\hypersetup{colorlinks=true, linkcolor=azulMedio, urlcolor=verde, citecolor=azul}

% =====================================================
% CABEÇALHO E RODAPÉ
% =====================================================
\pagestyle{fancy}
\fancyhf{}
\fancyhead[L]{%
    \IfFileExists{../figures/garopaba_horizontal_marca2015_PNG.png}{%
        \raisebox{-0.25\height}{\includegraphics[height=0.95cm]{../figures/garopaba_horizontal_marca2015_PNG.png}}%
    }{\textbf{\color{verde} IFSC — Câmpus Garopaba}}}
\fancyhead[R]{\footnotesize\color{suave}\textbf{\color{azul}Roteiro Prático \numeroRoteiro}\hspace{0.5em}\textbullet\hspace{0.5em}\tituloCurto}
\fancyfoot[L]{\footnotesize\color{suave}Programação para a Internet 2\hspace{0.4em}\textbullet\hspace{0.4em}CST Sistemas para a Internet}
\fancyfoot[R]{\tcbox[on line, colback=azul!8, colframe=azul!8, arc=6pt, boxrule=0pt, left=5pt, right=5pt, top=1pt, bottom=1pt, boxsep=0pt]{\footnotesize\color{azul}\bfseries\thepage\,/\,\pageref*{LastPage}}}
\renewcommand{\headrule}{{\color{borda}\hrule height 0.6pt}}
\renewcommand{\footrulewidth}{0pt}
\setlength{\headheight}{36pt}

% =====================================================
% TÍTULOS DE SEÇÃO
% =====================================================
\titleformat{\section}
  {\normalfont\Large\bfseries\color{azul}}
  {\tikz[baseline=(n.base)]\node[fill=azul, rounded corners=5pt, text=white, inner sep=5pt,
      minimum width=1.9em, font=\large\bfseries](n){\thesection};}
  {0.7em}{}
\titlespacing*{\section}{0pt}{1.6em}{0.7em}

\titleformat{\subsection}
  {\normalfont\large\bfseries\color{azul!85!black}}
  {{\color{verdeVivo}\thesubsection}}
  {0.6em}{}
\titlespacing*{\subsection}{0pt}{1.2em}{0.4em}

\titleformat{\subsubsection}
  {\normalfont\normalsize\bfseries\color{azul!85!black}}
  {{\color{verdeVivo}\thesubsubsection}}
  {0.6em}{}
\titlespacing*{\subsubsection}{0pt}{1.0em}{0.3em}

% Títulos nunca ficam sozinhos no pé da página
\pretocmd{\section}{\needspace{7\baselineskip}}{}{}
\pretocmd{\subsection}{\needspace{5\baselineskip}}{}{}

\setlist[itemize]{leftmargin=1.4em, itemsep=0.25em, topsep=0.3em, label={\color{verdeVivo}\small\faAngleRight}}
\setlist[enumerate]{leftmargin=2.2em, itemsep=0.45em, topsep=0.3em}

% Tabelas: sem linhas verticais, cabeçalho colorido, zebra suave
\newcommand{\cabecalho}{\rowcolor{azul}}
\newcommand{\cab}[1]{\leavevmode\color{white}\bfseries #1}
\arrayrulecolor{borda}
\newcolumntype{Y}{>{\raggedright\arraybackslash}X}
\newcolumntype{P}[1]{>{\raggedright\arraybackslash}p{#1}}
% Primeira coluna em destaque (rótulo da linha)
\newcolumntype{R}{>{\raggedright\arraybackslash\bfseries\leavevmode\color{azul}}l}
\newcolumntype{Q}[1]{>{\raggedright\arraybackslash\bfseries\leavevmode\color{azul}}p{#1}}

% Legendas de figuras e tabelas discretas
\usepackage[font={small,color=suave}, labelfont={bf,color=azul}, skip=6pt]{caption}

% \begin{document} ... (conteúdo do roteiro)

%   \begin{document}
%   \input{content.tex}
%   \end{document}
% Compilar com: tectonic main.tex
% =====================================================================
\documentclass[11pt,a4paper]{article}

\providecommand{\numeroRoteiro}{0X}
\providecommand{\tituloCurto}{Título curto do roteiro}

% =====================================================
% PACOTES
% =====================================================
\usepackage[utf8]{inputenc}
\usepackage[T1]{fontenc}
\usepackage[brazil]{babel}
\shorthandoff{"}
\AtBeginDocument{\shorthandoff{"}}
\usepackage[default]{sourcesanspro}          % texto: sem serifa, leitura confortável
\usepackage[scale=0.86]{sourcecodepro}       % código: monoespaçada moderna
% =====================================================
% ÍCONES (fontawesome5 com compatibilidade Tectonic no macOS)
% Evita crash de alocação na rotina nativa \XeTeXglyphname do Tectonic
% =====================================================
\usepackage{expl3}
\ExplSyntaxOn
\cs_set:Npn \tex_XeTeXglyphname:D #1#2 { }
\ExplSyntaxOff
\newcommand{\faDice}{}
\usepackage{fontawesome5}

\NewDocumentCommand\definirFaIcone{m m}{%
  \expandafter\NewDocumentCommand\csname fa#1\endcsname{O{}}{%
    \if\relax\detokenize{##1}\relax
      \faIcon{#2}%
    \else
      \faIcon[##1]{#2}%
    \fi
  }%
}
\definirFaIcone{AngleRight}{angle-right}
\definirFaIcone{BookOpen}{book-open}
\definirFaIcone{Bullseye}{bullseye}
\definirFaIcone{Calculator}{calculator}
\definirFaIcone{ChalkboardTeacher}{chalkboard-teacher}
\definirFaIcone{CheckCircle}{check-circle}
\definirFaIcone{ChevronRight}{chevron-right}
\definirFaIcone{Circle}{circle}
\definirFaIcone{ClipboardCheck}{clipboard-check}
\definirFaIcone{Code}{code}
\definirFaIcone{Columns}{columns}
\definirFaIcone{Database}{database}
\definirFaIcone{DraftingCompass}{drafting-compass}
\definirFaIcone{ExclamationTriangle}{exclamation-triangle}
\definirFaIcone{ExternalLinkAlt}{external-link-alt}
\definirFaIcone{Flag}{flag}
\definirFaIcone{FolderOpen}{folder-open}
\definirFaIcone{Github}{github}
\definirFaIcone{LaptopCode}{laptop-code}
\definirFaIcone{LayerGroup}{layer-group}
\definirFaIcone{Lightbulb}{lightbulb}
\definirFaIcone{Link}{link}
\definirFaIcone{Lock}{lock}
\definirFaIcone{Map}{map}
\definirFaIcone{MapMarkerAlt}{map-marker-alt}
\definirFaIcone{MousePointer}{mouse-pointer}
\definirFaIcone{Server}{server}
\definirFaIcone{Square}{square}
\definirFaIcone{Terminal}{terminal}
\definirFaIcone{Tools}{tools}
\definirFaIcone{University}{university}
\definirFaIcone{Check}{check}
\definirFaIcone{Times}{times}
\definirFaIcone{InfoCircle}{info-circle}
\definirFaIcone{QuestionCircle}{question-circle}
\definirFaIcone{Play}{play}
\definirFaIcone{FileCode}{file-code}
\definirFaIcone{FileAlt}{file-alt}
\definirFaIcone{Folder}{folder}
\definirFaIcone{Cogs}{cogs}
\definirFaIcone{Cog}{cog}
\definirFaIcone{Eye}{eye}
\definirFaIcone{Desktop}{desktop}
\definirFaIcone{Globe}{globe}
\definirFaIcone{Search}{search}
\definirFaIcone{Download}{download}
\definirFaIcone{Upload}{upload}
\definirFaIcone{Sync}{sync}
\definirFaIcone{ShareAlt}{share-alt}
\definirFaIcone{Edit}{edit}
\definirFaIcone{Trash}{trash}
\definirFaIcone{Plus}{plus}
\definirFaIcone{Minus}{minus}
\usepackage{geometry}
\usepackage{graphicx}
\graphicspath{{../figures/}{figures/}}
\usepackage{fancyhdr}
\usepackage{titlesec}
\usepackage[table]{xcolor}
\usepackage{hyperref}
\usepackage{xurl}                            % URLs longas quebram em qualquer ponto
\usepackage{enumitem}
\usepackage[most]{tcolorbox}
\usepackage{float}
\usepackage{tikz}
\usetikzlibrary{arrows.meta, positioning, shapes.geometric, calc, babel}
\usepackage{tabularx}
\usepackage{booktabs}
\usepackage{listings}
\usepackage{upquote}
\usepackage{lastpage}
\usepackage{qrcode}
\usepackage{multicol}
\usepackage{setspace}
\usepackage{needspace}
\usepackage{adjustbox}

\geometry{top=2.6cm, bottom=2.6cm, left=2cm, right=2cm, headsep=0.8cm}
\linespread{1.18}
\setlength{\parskip}{0.55em}
\setlength{\parindent}{0pt}
\setlength{\emergencystretch}{6em}   % evita linhas estouradas com \texttt longos

% =====================================================
% PALETA (cores institucionais IFSC, em tons suaves)
% =====================================================
\definecolor{azul}{RGB}{30,60,90}
\definecolor{azulMedio}{RGB}{37,99,235}
\definecolor{verde}{RGB}{22,101,52}
\definecolor{verdeVivo}{RGB}{22,163,74}
\definecolor{ambar}{RGB}{217,119,6}
\definecolor{vermelho}{RGB}{220,38,38}
\definecolor{roxo}{RGB}{124,58,237}
\definecolor{laranja}{RGB}{194,65,12}
\definecolor{ciano}{RGB}{14,116,144}
\definecolor{tinta}{RGB}{30,41,59}
\definecolor{suave}{RGB}{100,116,139}
\definecolor{borda}{RGB}{226,232,240}
\definecolor{fundo}{RGB}{248,250,252}
\definecolor{barraTitulo}{RGB}{241,245,249}
\definecolor{terminal}{RGB}{15,23,42}

% Nomes antigos dos roteiros de ProgWeb 2 (mantidos para compatibilidade)
\colorlet{azulTitulo}{azul}
\colorlet{azulObjetivos}{azul}
\colorlet{ifscGreen}{verdeVivo}
\colorlet{vermelhoCodigo}{vermelho}
\colorlet{cinzaCodigo}{fundo}
\colorlet{cinzaFundo}{fundo}
\colorlet{cinzaTabela}{fundo}
\colorlet{cinzaLinha}{fundo}
\colorlet{cinzaCabecalho}{barraTitulo}
\colorlet{pretoTabela}{borda}
\colorlet{borderGray}{borda}
\colorlet{bordaCodigo}{borda}
\colorlet{jsKeyword}{roxo}
\colorlet{jsString}{verde}
\colorlet{jsComment}{suave}
\colorlet{htmlTag}{azulMedio}
\colorlet{htmlString}{verde}
\colorlet{pyKeyword}{laranja}
\colorlet{pyBuiltin}{ciano}

\color{tinta}

% =====================================================
% LISTINGS — linguagens e estilos
% Regra de cópia: nenhuma linha de código pode passar de 103 caracteres
% (é o que cabe numa linha da janela de código; acima disso o PDF quebra a linha).
% =====================================================
\newcommand{\fonteCodigo}{\fontsize{8.5pt}{11.2pt}\selectfont\ttfamily}

\lstdefinelanguage{JavaScriptES6}{
  keywords={typeof, new, true, false, catch, function, return, null, switch, var, if, in, of, while, do, else, case, break, const, let, export, import, from, default, async, await, try, finally, throw, for, class, this},
  sensitive=true,
  comment=[l]{//},
  morecomment=[s]{/*}{*/},
  morestring=[b]',
  morestring=[b]",
  morestring=[b]`
}
\lstdefinelanguage{PythonCustom}{
  keywords={def, return, import, from, as, if, elif, else, while, for, in, try, except, finally, with, class, None, True, False, and, or, not, pass},
  ndkeywords={self, app, request, jsonify, CORS, sqlite3},
  sensitive=true,
  comment=[l]{\#},
  morestring=[b]',
  morestring=[b]",
  morestring=[s]{'''}{'''},
  morestring=[s]{"""}{"""}
}
\lstdefinelanguage{HTML5}{language=html}

\lstdefinestyle{base}{
    basicstyle=\fonteCodigo\color{tinta},
    keywordstyle=\color{roxo}\bfseries,
    ndkeywordstyle=\color{ciano}\bfseries,
    commentstyle=\color{suave}\itshape,
    stringstyle=\color{verde},
    identifierstyle=\color{tinta},
    breaklines=true, columns=fullflexible, keepspaces=true,
    showstringspaces=false, upquote=true, tabsize=2,
    frame=none, numbers=none, aboveskip=0pt, belowskip=0pt,
    extendedchars=true,
    literate=
      {á}{{\'a}}1 {é}{{\'e}}1 {í}{{\'i}}1 {ó}{{\'o}}1 {ú}{{\'u}}1
      {Á}{{\'A}}1 {É}{{\'E}}1 {Í}{{\'I}}1 {Ó}{{\'O}}1 {Ú}{{\'U}}1
      {à}{{\`a}}1 {è}{{\`e}}1 {ì}{{\`i}}1 {ò}{{\`o}}1 {ù}{{\`u}}1
      {ã}{{\~a}}1 {õ}{{\~o}}1 {Ã}{{\~A}}1 {Õ}{{\~O}}1
      {â}{{\^a}}1 {ê}{{\^e}}1 {î}{{\^i}}1 {ô}{{\^o}}1 {û}{{\^u}}1
      {Â}{{\^A}}1 {Ê}{{\^E}}1 {Ô}{{\^O}}1
      {ç}{{\c c}}1 {Ç}{{\c C}}1
      {–}{{--}}1 {—}{{---}}1 {±}{{$\pm$}}1 {≈}{{$\approx$}}1 {÷}{{$\div$}}1
}
\lstdefinestyle{js}{style=base, language=JavaScriptES6}
\lstdefinestyle{py}{style=base, language=PythonCustom, keywordstyle=\color{laranja}\bfseries, tabsize=4}
\lstdefinestyle{html}{style=base, language=HTML5, tagstyle=\color{azulMedio}\bfseries,
    keywordstyle=\color{azulMedio}, stringstyle=\color{verde}}
\lstdefinestyle{txt}{style=base, language={}}
\lstdefinestyle{saida}{style=base, language={}, basicstyle=\fonteCodigo\color{white!88!black},
    identifierstyle=\color{white!88!black}}
\lstset{style=base}

% =====================================================
% COMPONENTES VISUAIS
% =====================================================
% --- Janela de código (estilo editor) ---------------
% \begin{codigo}[js|html|py|txt]{nome-do-arquivo} ... \end{codigo}
\newcommand{\pontosJanela}{%
  {\color{vermelho!70}\tiny\faCircle}\,{\color{ambar!80}\tiny\faCircle}\,{\color{verdeVivo!70}\tiny\faCircle}\hspace{0.8em}}
\tcbset{janela/.style={
    enhanced, listing only, breakable,
    colback=fundo, colframe=borda, boxrule=0.6pt, arc=6pt,
    left=8pt, right=8pt, top=6pt, bottom=6pt,
    coltitle=suave, colbacktitle=barraTitulo,
    fonttitle=\ttfamily\footnotesize, toptitle=2pt, bottomtitle=2pt,
    titlerule=0.6pt, titlerule style=borda,
    before skip=6pt, after skip=10pt}}
% Etiqueta da linguagem, à direita da barra de título
\newcommand{\rotuloLinguagem}[2]{\hfill{\sffamily\scriptsize\bfseries\color{#1}#2}}
\newtcblisting{codigo}[2][js]{janela, title={\pontosJanela #2}, listing options={style=#1}}

% --- Janelas por linguagem (nomes usados nos roteiros) ----------------
% O argumento opcional é o nome do arquivo mostrado na barra da janela.
\newtcblisting{boxCodigoJS}[1][]{janela,
    title={\pontosJanela#1\rotuloLinguagem{ambar}{JavaScript}},
    listing options={style=js}}
\newtcblisting{boxCodigoHTML}[1][]{janela,
    title={\pontosJanela#1\rotuloLinguagem{vermelho}{HTML}},
    listing options={style=html}}
\newtcblisting{boxCodigoPY}[1][]{janela,
    title={\pontosJanela#1\rotuloLinguagem{azulMedio}{Python}},
    listing options={style=py}}
\newtcblisting{boxCodigoCSS}[1][]{janela,
    title={\pontosJanela#1\rotuloLinguagem{roxo}{CSS}},
    listing options={style=txt}}
% Croqui (desenho em texto da tela final)
\newtcblisting{boxCroqui}[1][]{janela,
    title={{\color{azulMedio}\faDraftingCompass}\hspace{0.5em}{\sffamily\bfseries #1}},
    listing options={style=txt}}

% Árvore de pastas do projeto
\newtcblisting{arvore}[1][Estrutura de pastas do projeto]{janela,
    title={{\color{ambar}\faFolderOpen}\hspace{0.5em}{\sffamily\bfseries #1}},
    listing options={style=txt}}

% --- Saída / terminal (fundo escuro) -------------------------------
\newtcblisting{saida}[1][Terminal]{
    enhanced, listing only, breakable,
    colback=terminal, colframe=terminal, boxrule=0pt, arc=6pt,
    left=10pt, right=10pt, top=4pt, bottom=6pt,
    title={\faTerminal\hspace{0.5em}#1}, coltitle=white!65!black,
    colbacktitle=terminal, fonttitle=\sffamily\footnotesize\bfseries,
    toptitle=3pt, bottomtitle=0pt, titlerule=0pt,
    before skip=6pt, after skip=10pt,
    listing options={style=saida}
}
% Versões antigas (conteúdo LaTeX com verbatim dentro)
\newtcolorbox{boxCodigoTerminal}{
    enhanced, breakable,
    colback=terminal, colframe=terminal, boxrule=0pt, arc=6pt, coltext=white!88!black,
    left=10pt, right=10pt, top=4pt, bottom=6pt,
    title={\faTerminal\hspace{0.5em}Terminal}, coltitle=white!65!black,
    colbacktitle=terminal, fonttitle=\sffamily\footnotesize\bfseries,
    toptitle=3pt, bottomtitle=0pt, titlerule=0pt,
    before upper={\singlespacing}
}
\newtcolorbox{boxCodigo}{
    enhanced, breakable,
    colback=fundo, colframe=borda, boxrule=0.6pt, arc=6pt,
    left=10pt, right=10pt, top=8pt, bottom=8pt,
    before upper={\singlespacing}
}

% --- Exemplo certo/errado -----------------------------
\newtcblisting{exemplo}[3][js]{
    enhanced, listing only,
    colback=#2!4, colframe=#2!45, boxrule=0.6pt, arc=6pt,
    left=8pt, right=8pt, top=5pt, bottom=5pt,
    title={#3}, colbacktitle=#2!12, coltitle=#2!75!black,
    fonttitle=\sffamily\small\bfseries, toptitle=2pt, bottomtitle=2pt,
    listing options={style=#1}
}

% --- Caixas de destaque com barra lateral -----------
\newtcolorbox{destaque}[3]{
    enhanced, breakable,
    colback=#1!6, frame hidden, arc=6pt,
    borderline west={3.5pt}{0pt}{#1},
    left=14pt, right=12pt, top=9pt, bottom=9pt,
    before skip=8pt, after skip=10pt,
    before upper={{\color{#1}\large #2}\hspace{0.55em}{\color{#1!75!black}\bfseries #3}\par\vspace{0.25em}}
}
\newenvironment{dica}[1]{\begin{destaque}{ambar}{\faLightbulb[regular]}{#1}}{\end{destaque}}
\newenvironment{atencao}[1]{\begin{destaque}{vermelho}{\faExclamationTriangle}{#1}}{\end{destaque}}
\newenvironment{conceito}[1]{\begin{destaque}{azulMedio}{\faBookOpen}{#1}}{\end{destaque}}
\newenvironment{sucesso}[1]{\begin{destaque}{verdeVivo}{\faCheckCircle[regular]}{#1}}{\end{destaque}}
\newenvironment{avancado}[1]{\begin{destaque}{roxo}{\faCode}{#1}}{\end{destaque}}
\newtcolorbox{repositorio}[1]{enhanced, colback=azul!5, frame hidden, arc=6pt,
    borderline west={3.5pt}{0pt}{azul}, left=14pt, right=12pt, top=9pt, bottom=9pt,
    before skip=8pt, after skip=10pt, fontupper=\small,
    before upper={{\color{azul}\large\faGithub}\hspace{0.55em}{\color{azul!75!black}\bfseries #1}\par\vspace{0.25em}}}
\newenvironment{checkpoint}[1]{\begin{destaque}{verdeVivo}{\faClipboardCheck}{#1}}{\end{destaque}}

% Nomes antigos
\newenvironment{boxDica}[1][Dica]{\begin{dica}{#1}}{\end{dica}}
\newenvironment{boxTeste}[1][Checkpoint de validação]{\begin{checkpoint}{#1}}{\end{checkpoint}}

% --- Caixa com título "aba" (objetivos, atividade) ---
\newtcolorbox{painel}[3][]{
    enhanced, breakable,
    colback=white, colframe=borda, boxrule=0.6pt, arc=8pt,
    left=14pt, right=14pt, top=14pt, bottom=10pt,
    title={#2\hspace{0.5em}#3}, coltitle=white, colbacktitle=azul,
    fonttitle=\bfseries, attach boxed title to top left={xshift=12pt, yshift=-8pt},
    boxed title style={arc=4pt, boxrule=0pt},
    before skip=12pt, after skip=10pt, #1
}
\newenvironment{boxObjetivos}[1][Ao final desta prática você será capaz de\ldots]
  {\begin{painel}{\faBullseye}{#1}}{\end{painel}}
\newenvironment{boxAtividade}[1][Atividade]
  {\begin{painel}[colbacktitle=azulMedio]{\faLaptopCode}{#1}}{\end{painel}}
\newtcolorbox{boxTabela}{
    enhanced, colback=white, colframe=borda, boxrule=0.6pt, arc=6pt,
    left=6pt, right=6pt, top=6pt, bottom=6pt}

% --- Checklist de entrega (itens com \caixaMarcar) ----------------
\newtcolorbox{checklist}[1][Antes de entregar, confira se\ldots]{enhanced, breakable,
    colback=verdeVivo!5, colframe=verdeVivo!40, boxrule=0.6pt, arc=8pt,
    left=14pt, right=14pt, top=10pt, bottom=10pt,
    title={\faClipboardCheck\hspace{0.5em}#1}, colbacktitle=verdeVivo, coltitle=white, fonttitle=\bfseries,
    before skip=10pt, after skip=10pt}

% --- Navegador (prévia da tela) -----------------------
\newtcolorbox{mockupNavegador}[1][]{
    enhanced, breakable,
    colback=white, colframe=borda, boxrule=0.8pt, arc=7pt,
    left=8pt, right=8pt, top=8pt, bottom=8pt,
    title={\pontosJanela\tcbox[on line, colback=white, colframe=borda, boxrule=0.5pt, arc=6pt,
        left=6pt, right=6pt, top=0.5pt, bottom=0.5pt, boxsep=0.5pt]{\ttfamily\scriptsize\color{suave}\faLock\ 127.0.0.1:5500/index.html}%
        \hfill{\sffamily\footnotesize\bfseries\color{suave}#1}},
    colbacktitle=barraTitulo, coltitle=tinta, titlerule=0.6pt, titlerule style=borda,
    toptitle=2pt, bottomtitle=2pt
}

% --- Cartões pequenos ---------------------------------
\tcbset{cartao/.style={enhanced, colback=fundo, colframe=borda, boxrule=0.6pt, arc=8pt,
    left=10pt, right=10pt, top=8pt, bottom=8pt, fontupper=\small\raggedright}}

% --- Elementos em linha --------------------------------
\newcommand{\tecla}[1]{\tcbox[on line, enhanced, colback=white, colframe=suave!50, boxrule=0.5pt,
    arc=2.5pt, left=3pt, right=3pt, top=0.5pt, bottom=0.5pt, boxsep=0.5pt,
    borderline south={1.2pt}{0pt}{suave!40}]{\ttfamily\footnotesize #1}}
\newcommand{\botao}[1]{\tcbox[on line, enhanced, colback=azulMedio!8, colframe=azulMedio!35, boxrule=0.5pt,
    arc=3pt, left=3pt, right=3pt, top=0.5pt, bottom=0.5pt, boxsep=0.5pt]{\sffamily\footnotesize\color{azul}#1}}
\newcommand{\seta}{\,{\color{suave}\scriptsize\faChevronRight}\,}
\newcommand{\chip}[2]{\tcbox[on line, enhanced, colback=white, opacityback=0.14, colframe=white,
    opacityframe=0.45, boxrule=0.5pt, arc=8pt, left=5pt, right=5pt, top=1.5pt, bottom=1.5pt,
    boxsep=0pt, coltext=white]{\footnotesize #1\hspace{0.4em}#2}}
\newcommand{\circulo}[2][azul]{\tikz[baseline=(c.base)]\node[circle, fill=#1, text=white,
    inner sep=0pt, minimum size=1.35em, font=\footnotesize\bfseries](c){#2};}
\newcommand{\marca}[1]{\tikz[baseline=(c.base)]\node[circle, fill=vermelho, text=white,
    inner sep=0pt, minimum size=1.2em, font=\scriptsize\bfseries](c){#1};}
\newcommand{\caixaMarcar}{{\color{suave}\faSquare[regular]}}

% --- Mantém um bloco de código curto inteiro na mesma página --------
% \manterJunto{n}: n = número de linhas do código que vem a seguir
\newcommand{\manterJunto}[1]{\needspace{\dimexpr 11.2pt*#1+5\baselineskip\relax}}

% --- Passo numerado do passo a passo ------------------
% \passo{1.2}{Criar a página HTML}  — o nome do arquivo vai na barra da janela de código
\newcommand{\passo}[2]{%
  \par\vspace{0.5em}\noindent
  \tcbox[on line, colback=azul, colframe=azul, arc=4pt, boxrule=0pt,
      left=5pt, right=5pt, top=1.5pt, bottom=1.5pt, boxsep=0pt]{\sffamily\footnotesize\bfseries\color{white}Passo #1}%
  \hspace{0.55em}{\bfseries\color{azul}#2}%
  \par\nopagebreak\vspace{-0.2em}}

% =====================================================
% ABERTURA PADRÃO DO ROTEIRO
% =====================================================
% \capaRoteiro{ícone-marca-d'água}{Título}{Subtítulo}{chips}
\newcommand{\capaRoteiro}[4]{%
\begin{tcolorbox}[enhanced, frame hidden, arc=12pt,
    interior style={left color=azul, right color=verde},
    left=20pt, right=20pt, top=18pt, bottom=16pt, colupper=white,
    overlay={\node[anchor=east, opacity=0.13, font=\fontsize{90}{90}\selectfont, text=white] at ([xshift=-14pt]frame.east) {#1};}]
{\footnotesize\bfseries\color{white!80!azul}ROTEIRO PRÁTICO \numeroRoteiro\hspace{0.6em}\textbullet\hspace{0.6em}PROGRAMAÇÃO PARA A INTERNET 2}\par\vspace{2pt}
{\fontsize{24}{29}\selectfont\bfseries #2}\par\vspace{2pt}
{\large #3}\par\vspace{10pt}
#4
\end{tcolorbox}}

% \cartoesUC{Produto/entrega deste roteiro}
\newcommand{\cartoesUC}[1]{%
\begin{tcbraster}[raster columns=3, raster equal height, raster column skip=8pt]
\begin{tcolorbox}[cartao]
{\color{azulMedio}\faUniversity}\hspace{0.4em}\textbf{Unidade Curricular}\par\vspace{2pt}
Programação para a Internet 2\\ {\color{suave}CST em Sistemas para a Internet}
\end{tcolorbox}
\begin{tcolorbox}[cartao]
{\color{azulMedio}\faChalkboardTeacher}\hspace{0.4em}\textbf{Docente}\par\vspace{2pt}
Prof. André Moraes\\ {\color{suave}IFSC Câmpus Garopaba}
\end{tcolorbox}
\begin{tcolorbox}[cartao]
{\color{azulMedio}\faFlag}\hspace{0.4em}\textbf{Produto deste roteiro}\par\vspace{2pt}
#1
\end{tcolorbox}
\end{tcbraster}}

% \etapaMapa{cor}{n}{Título}{descrição}  — usar dentro de \begin{mapaAula} ... \end{mapaAula}
\newenvironment{mapaAula}{%
  \needspace{9\baselineskip}{\large\bfseries\color{azul}\faMap[regular]\hspace{0.5em}Mapa da aula}\par\vspace{0.2em}
  \begin{tcbraster}[raster columns=2, raster equal height, raster column skip=8pt, raster row skip=6pt]}
  {\end{tcbraster}}
\newcommand{\etapaMapa}[4]{%
\begin{tcolorbox}[cartao, colback=#1!5, colframe=#1!25]
\circulo[#1]{#2}\hspace{0.4em}\textbf{#3}\par\vspace{2pt}
{\color{suave}\footnotesize #4}
\end{tcolorbox}}

% \trilhaRoteiros{n} — trilha da disciplina com o roteiro atual em destaque
\newcommand{\trilhaRoteiros}[1]{%
\begin{center}
\begin{tikzpicture}[node distance=0.32cm,
    etapa/.style={rectangle, draw=borda, fill=fundo, line width=0.6pt, text width=2.55cm, align=center,
        rounded corners=7pt, minimum height=2.0cm, font=\footnotesize},
    atual/.style={draw=verdeVivo, fill=verdeVivo!7, line width=1.1pt},
    linha/.style={draw=suave!70, -{Latex[length=2.2mm]}, line width=0.9pt},
    r1/.style={}, r2/.style={}, r3/.style={}, r4/.style={}, r5/.style={}, r#1/.style={atual}]
    \node[etapa, r1] (r1) {{\color{azulMedio}\Large\faColumns}\\[3pt]\textbf{Roteiro 01}\\{\color{suave}Bootstrap 5 e layout}};
    \node[etapa, right=of r1, r2] (r2) {{\color{azulMedio}\Large\faMousePointer}\\[3pt]\textbf{Roteiro 02}\\{\color{suave}jQuery e DOM}};
    \node[etapa, right=of r2, r3] (r3) {{\color{azulMedio}\Large\faCalculator}\\[3pt]\textbf{Roteiro 03}\\{\color{suave}Calculadora e LocalStorage}};
    \node[etapa, right=of r3, r4] (r4) {{\color{azulMedio}\Large\faLayerGroup}\\[3pt]\textbf{Roteiro 04}\\{\color{suave}ES6 Modules em camadas}};
    \node[etapa, right=of r4, r5] (r5) {{\color{azulMedio}\Large\faServer}\\[3pt]\textbf{Roteiro 05}\\{\color{suave}Fetch API e Flask}};
    \draw[linha] (r1) -- (r2);
    \draw[linha] (r2) -- (r3);
    \draw[linha] (r3) -- (r4);
    \draw[linha] (r4) -- (r5);
\end{tikzpicture}
\end{center}}

% =====================================================
% HIPERLINKS
% =====================================================
\hypersetup{colorlinks=true, linkcolor=azulMedio, urlcolor=verde, citecolor=azul}

% =====================================================
% CABEÇALHO E RODAPÉ
% =====================================================
\pagestyle{fancy}
\fancyhf{}
\fancyhead[L]{%
    \IfFileExists{../figures/garopaba_horizontal_marca2015_PNG.png}{%
        \raisebox{-0.25\height}{\includegraphics[height=0.95cm]{../figures/garopaba_horizontal_marca2015_PNG.png}}%
    }{\textbf{\color{verde} IFSC — Câmpus Garopaba}}}
\fancyhead[R]{\footnotesize\color{suave}\textbf{\color{azul}Roteiro Prático \numeroRoteiro}\hspace{0.5em}\textbullet\hspace{0.5em}\tituloCurto}
\fancyfoot[L]{\footnotesize\color{suave}Programação para a Internet 2\hspace{0.4em}\textbullet\hspace{0.4em}CST Sistemas para a Internet}
\fancyfoot[R]{\tcbox[on line, colback=azul!8, colframe=azul!8, arc=6pt, boxrule=0pt, left=5pt, right=5pt, top=1pt, bottom=1pt, boxsep=0pt]{\footnotesize\color{azul}\bfseries\thepage\,/\,\pageref*{LastPage}}}
\renewcommand{\headrule}{{\color{borda}\hrule height 0.6pt}}
\renewcommand{\footrulewidth}{0pt}
\setlength{\headheight}{36pt}

% =====================================================
% TÍTULOS DE SEÇÃO
% =====================================================
\titleformat{\section}
  {\normalfont\Large\bfseries\color{azul}}
  {\tikz[baseline=(n.base)]\node[fill=azul, rounded corners=5pt, text=white, inner sep=5pt,
      minimum width=1.9em, font=\large\bfseries](n){\thesection};}
  {0.7em}{}
\titlespacing*{\section}{0pt}{1.6em}{0.7em}

\titleformat{\subsection}
  {\normalfont\large\bfseries\color{azul!85!black}}
  {{\color{verdeVivo}\thesubsection}}
  {0.6em}{}
\titlespacing*{\subsection}{0pt}{1.2em}{0.4em}

\titleformat{\subsubsection}
  {\normalfont\normalsize\bfseries\color{azul!85!black}}
  {{\color{verdeVivo}\thesubsubsection}}
  {0.6em}{}
\titlespacing*{\subsubsection}{0pt}{1.0em}{0.3em}

% Títulos nunca ficam sozinhos no pé da página
\pretocmd{\section}{\needspace{7\baselineskip}}{}{}
\pretocmd{\subsection}{\needspace{5\baselineskip}}{}{}

\setlist[itemize]{leftmargin=1.4em, itemsep=0.25em, topsep=0.3em, label={\color{verdeVivo}\small\faAngleRight}}
\setlist[enumerate]{leftmargin=2.2em, itemsep=0.45em, topsep=0.3em}

% Tabelas: sem linhas verticais, cabeçalho colorido, zebra suave
\newcommand{\cabecalho}{\rowcolor{azul}}
\newcommand{\cab}[1]{\leavevmode\color{white}\bfseries #1}
\arrayrulecolor{borda}
\newcolumntype{Y}{>{\raggedright\arraybackslash}X}
\newcolumntype{P}[1]{>{\raggedright\arraybackslash}p{#1}}
% Primeira coluna em destaque (rótulo da linha)
\newcolumntype{R}{>{\raggedright\arraybackslash\bfseries\leavevmode\color{azul}}l}
\newcolumntype{Q}[1]{>{\raggedright\arraybackslash\bfseries\leavevmode\color{azul}}p{#1}}

% Legendas de figuras e tabelas discretas
\usepackage[font={small,color=suave}, labelfont={bf,color=azul}, skip=6pt]{caption}

% \begin{document} ... (conteúdo do roteiro)


\begin{document}

% =====================================================
% CAPA / ABERTURA INSTITUCIONAL
% =====================================================
\capaRoteiro{\faCogs}{Arquitetura Modular do Frontend}
  {Documentação técnica da estrutura em camadas com ES6 Modules e desacoplamento do DOM}
  {\chip{\faLayerGroup}{Arquitetura Web}\hspace{4pt}%
   \chip{\faCode}{ES6 Modules}\hspace{4pt}%
   \chip{\faDatabase}{Storage \& API}\hspace{4pt}%
   \chip{\faUniversity}{IFSC Garopaba}}

\vspace{0.4em}
\cartoesUC{Registro da arquitetura modular do portal, guia de manutenção e referência pedagógica viva}

\begin{boxObjetivos}
\begin{itemize}
  \item \textbf{Registrar formalmente a arquitetura modular} adotada no portal didático da disciplina de Programação para a Internet 2;
  \item \textbf{Descrever a divisão em quatro camadas coesas} (\texttt{config/}, \texttt{utils/}, \texttt{services/} e \texttt{views/}), alinhadas aos conteúdos ensinados no Roteiro 04;
  \item \textbf{Documentar a eliminação do monólito de 5.445 linhas de script inline} do \texttt{index.html}, demonstrando ganho de performance e manutenibilidade;
  \item \textbf{Explicar a ponte de retrocompatibilidade global} estabelecida via \texttt{app.js} para elementos declarativos do DOM legados;
  \item \textbf{Fornecer um guia cirúrgico para novas funcionalidades}, permitindo adicionar demandas, serviços e rotas com isolamento e segurança;
  \item \textbf{Apresentar a matriz de métricas e testes de regressão} que asseguram a estabilidade do portal no GitHub Pages.
\end{itemize}
\end{boxObjetivos}

\begin{mapaAula}
\etapaMapa{azulMedio}{1}{Contexto \& Motivação}{Gargalos do monólito e alinhamento pedagógico (Seção~\ref{sec:contexto})}
\etapaMapa{verdeVivo}{2}{Visão Geral das Camadas}{Diagrama de fluxo, contratos e responsabilidades (Seção~\ref{sec:camadas})}
\etapaMapa{ambar}{3}{Estrutura de Arquivos}{Módulos em \texttt{js/}, serviços e ponte com o DOM (Seção~\ref{sec:estrutura})}
\etapaMapa{roxo}{4}{Guia de Extensão \& Métricas}{Como estender o portal e matriz de impacto (Seção~\ref{sec:extensao})}
\end{mapaAula}

% =====================================================
% 1. CONTEXTO E MOTIVAÇÃO DA REFATORAÇÃO
% =====================================================
\section{Contexto e Motivação da Refatoração}\label{sec:contexto}

Durante o desenvolvimento inicial do portal unificado da disciplina de Programação para a Internet 2, o arquivo \path{index.html} concentrava tanto a marcação semântica e os modais Bootstrap 5 quanto a totalidade da lógica em JavaScript da aplicação. Tratava-se de um monólito com \textbf{mais de 5.445 linhas de código inline} contidas em uma única tag \texttt{<script>}, elevando o tamanho do documento para impressionantes 11.070 linhas.

\begin{conceito}{A Importância do Desacoplamento Pedagógico}
Em sala de aula, ensina-se aos futuros tecnólogos a separação estrita de interesses (\textit{Separation of Concerns}). A prática pedagógica é consolidada no \textbf{Roteiro Prático 04 — Arquitetura Modular com ES6 Modules}, no qual os estudantes estruturam aplicações web separando utilitários, serviços de persistência e camadas de visualização. Manter o portal da própria unidade curricular em formato monolítico criava uma contradição didática que a presente refatoração eliminou integralmente.
\end{conceito}

Esse cenário anterior gerava três ordens de desafios críticos:
\begin{enumerate}
  \item \textbf{Sobrecarga de Contexto e Degradação de Manutenção:} A adição ou correção de qualquer elemento visual (como filtros da Central de Demandas ou cartões da SNCT) exigia a inspeção de milhares de linhas, aumentando a latência das intervenções e o risco de colisões;
  \item \textbf{Acoplamento Severo com o DOM e Armazenamento:} Constantes globais, manipulação direta de \texttt{localStorage}, rotinas de OAuth2 com Google Identity Services e geradores de templates HTML compartilhavam o mesmo escopo global sem limites arquiteturais;
  \item \textbf{Dificuldade na Adição de Testes e Evolução:} Funções utilitárias como sanitização de strings contra XSS (\texttt{escapeHtml}) e formatação temporal ficavam inacessíveis para reuso isolado.
\end{enumerate}

Com a refatoração concluída em 4 fases cirúrgicas, o \path{index.html} foi reduzido em \textbf{5.448 linhas (corte de 49,2\%)}, tornando-se um arquivo estritamente declarativo, enquanto toda a inteligência da aplicação passou a operar via módulos nativos do ES6 (\path{<script type="module" src="js/app.js">}).

% =====================================================
% 2. VISÃO GERAL DA ARQUITETURA EM CAMADAS
% =====================================================
\section{Visão Geral da Arquitetura em Camadas}\label{sec:camadas}

A aplicação foi estruturada seguindo o padrão de camadas unidirecionais, onde a apresentação depende dos serviços e configurações, mas as camadas inferiores jamais conhecem detalhes de interface:

\begin{center}
\begin{adjustbox}{max width=\linewidth}
\begin{tikzpicture}[
    node distance=1.1cm and 1.8cm,
    caixaDOM/.style={rectangle, draw=azulTitulo, rounded corners=6pt, fill=azulTitulo!8, line width=1.1pt, align=center, inner sep=7pt, font=\small},
    caixaApp/.style={rectangle, draw=ambar, rounded corners=6pt, fill=ambar!10, line width=1.1pt, align=center, inner sep=7pt, font=\small\bfseries},
    caixaView/.style={rectangle, draw=ifscGreen, rounded corners=6pt, fill=ifscGreen!8, line width=1.1pt, align=center, inner sep=6pt, font=\small},
    caixaService/.style={rectangle, draw=roxo, rounded corners=6pt, fill=roxo!8, line width=1.1pt, align=center, inner sep=6pt, font=\small},
    caixaBase/.style={rectangle, draw=suave, rounded corners=6pt, fill=fundo, line width=1pt, align=center, inner sep=6pt, font=\footnotesize},
    seta/.style={-{Latex[length=2.8mm]}, line width=1.1pt}
]

% Linha 1: DOM e App
\node[caixaDOM] (dom) {\textbf{index.html (DOM)}\\Marcação Bootstrap 5 e Modais};
\node[caixaApp, right=of dom] (app) {js/app.js (Orquestrador)\\Estado Reativo e Ponte Global};

% Linha 2: Camada de Views
\node[caixaView, below left=1.2cm and -0.8cm of app] (vDem) {\textbf{views/demandasView.js}\\Cards, Prazos e Badges};
\node[caixaView, right=0.4cm of vDem] (vSnct) {\textbf{views/snctView.js}\\5 Etapas SNCT e Vagas};
\node[caixaView, right=0.4cm of vSnct] (vRub) {\textbf{views/rubricasView.js}\\Modal de Rubricas (A2/A3/A4)};
\node[caixaView, right=0.4cm of vRub] (vNav) {\textbf{views/navigationView.js}\\Abas e Temas};

% Linha 3: Camada de Services
\node[caixaService, below=1.2cm of vSnct] (sStore) {\textbf{services/storageService.js}\\LocalStorage com Fallback};
\node[caixaService, left=0.6cm of sStore] (sApi) {\textbf{services/apiService.js}\\Fetch API e Google Apps Script};
\node[caixaService, right=0.6cm of sStore] (sAuth) {\textbf{services/authService.js}\\Google Identity Services};

% Linha 4: Config e Utils
\node[caixaBase, below=1.0cm of sStore] (conf) {\textbf{js/config/} (\texttt{constants.js}, \texttt{rubricasData.js}) \quad | \quad \textbf{js/utils/} (\texttt{formatters.js}, \texttt{toast.js})};

% Conexões
\draw[seta, color=azulTitulo] (dom) -- node[midway, above, font=\footnotesize] {Eventos / DOM} (app);
\draw[seta, color=ambar] (app.south) -- (vDem.north);
\draw[seta, color=ambar] (app.south) -- (vSnct.north);
\draw[seta, color=ambar] (app.south) -- (vRub.north);
\draw[seta, color=ambar] (app.south) -- (vNav.north);

\draw[seta, color=ifscGreen] (vDem.south) -- (sStore.north);
\draw[seta, color=ifscGreen] (vSnct.south) -- (sStore.north);
\draw[seta, color=ifscGreen] (vRub.south) -- (sStore.north);

\draw[seta, color=roxo] (sStore.south) -- (conf.north);
\draw[seta, color=roxo] (sApi.south) -- (conf.north);
\draw[seta, color=roxo] (sAuth.south) -- (conf.north);

\end{tikzpicture}
\end{adjustbox}
\end{center}

\begin{table}[H]
\centering
\small
\renewcommand{\arraystretch}{1.45}\rowcolors{2}{fundo}{white}
\begin{tabularx}{\textwidth}{RYY}
\cabecalho \cab{Camada} & \cab{Responsabilidade Arquitetural} & \cab{Principais Módulos} \\
\textbf{Apresentação (\texttt{views/})} & Geração declarativa de HTML dinâmico, atualização de badges, gerenciamento de modais e captura de ações do usuário. & \path{demandasView.js}, \path{snctView.js}, \path{rubricasView.js}, \path{navigationView.js} \\
\textbf{Serviços (\texttt{services/})} & Isolamento de operações de entrada e saída (I/O), persistência local, chamadas HTTP remotas e autenticação OAuth2. & \path{storageService.js}, \path{apiService.js}, \path{authService.js} \\
\textbf{Utilitários (\texttt{utils/})} & Funções puras sem efeitos colaterais para sanitização contra XSS, conversão de datas e toasts flutuantes. & \path{formatters.js}, \path{toast.js} \\
\textbf{Configuração (\texttt{config/})} & Centralização de chaves de storage, roster dos 24 alunos, frentes da SNCT e matrizes pedagógicas formativas. & \path{constants.js}, \path{rubricasData.js} \\
\textbf{Orquestrador (\texttt{app.js})} & Ponto de entrada modular, gerenciador do estado em memória e ponte de retrocompatibilidade global para o HTML. & \path{app.js}, \path{main.js} \\
\end{tabularx}
\caption{Matriz de responsabilidades das camadas da aplicação.}
\end{table}

% =====================================================
% 3. ESTRUTURA DE DIRETÓRIOS E CONTRATOS
% =====================================================
\section{Estrutura de Diretórios e Contratos de Módulos}\label{sec:estrutura}

A árvore de arquivos do portal reflete de forma fidedigna as diretrizes de engenharia de software frontend modernas:

\begin{arvore}[Estrutura Modular do Portal da UC ProgWeb 2]
uc-programacao-para-a-internet-2/
|-- index.html                 # Dashboard da disciplina (sem blocos <script> inline)
|-- progweb2-config.js         # Parametrização global de ambiente (URLs e Client ID)
`-- js/
    |-- app.js                 # Ponto de entrada carregado via <script type="module">
    |-- main.js                # Arquivo de barramento e re-exportação de contratos
    |-- config/
    |   |-- constants.js       # Chaves de armazenamento, roster de 24 alunos e frentes SNCT
    |   `-- rubricasData.js    # Matrizes contextuais de rubricas formativas (A2, A3 e A4)
    |-- services/
    |   |-- apiService.js      # Integração assíncrona com backend Google Apps Script
    |   |-- authService.js     # Autenticação Google GIS, decodificação JWT e vinculação
    |   `-- storageService.js  # Operações de CRUD com LocalStorage e fallback tolerante
    |-- utils/
    |   |-- formatters.js      # Sanitização XSS, formatação de datas e cálculo de prazos
    |   `-- toast.js           # Notificações visuais flutuantes com Bootstrap Toasts
    `-- views/
        |-- demandasView.js    # Renderização da Central de Demandas, cards e filtros
        |-- navigationView.js  # Chaveamento de abas, tema claro/escuro e foco de card
        |-- rubricasView.js    # Painel interativo de rubricas e gerador de parecer
        `-- snctView.js        # Diagrama das 5 etapas da SNCT e quadro de vagas
\end{arvore}

\subsection{Camada de Configuração (\texttt{js/config/})}

O módulo \path{constants.js} concentra todos os literais imutáveis do portal. Nenhuma chave de armazenamento é escrita como texto solto no código-fonte:

\begin{boxCodigoJS}[js/config/constants.js]
export const STORAGE_KEY_DEMANDAS = "progweb2_demandas_v1";
export const STORAGE_KEY_AVALIACOES = "progweb2_avaliacoes_v1";
export const STORAGE_KEY_ALUNOS = "progweb2_alunos_v1";
export const STORAGE_KEY_SNCT_DEMANDAS = "progweb2_snct_demandas_v1";
export const STORAGE_KEY_GOOGLE_SESSION = "progweb2_google_session_v1";

// Roster oficial completo com os 24 estudantes da turma 2026/2
export const ESTUDANTES_PADRAO = [
  { id: "aluno_1", nome: "Alexandre Silva", matricula: "2024101001", email: "alexandre@aluno.ifsc.edu.br" },
  { id: "aluno_2", nome: "Beatriz Oliveira", matricula: "2024101002", email: "beatriz@aluno.ifsc.edu.br" },
  // ... 22 demais discentes devidamente catalogados
];
\end{boxCodigoJS}

Já o arquivo \path{rubricasData.js} formaliza as três matrizes pedagógicas adotadas nas avaliações formativas da unidade curricular:
\begin{itemize}
  \item \textbf{Matriz A2 (SNCT):} Focada em planejamento didático, rigor conceitual da oficina, ambientação na sala Moodle e trabalho em equipe;
  \item \textbf{Matriz A3 (Projeto Integrador):} Focada em backend RESTful com Flask/SQLite, integração com Fetch API, contingência offline e ergonomia visual;
  \item \textbf{Matriz A4 (Roteiros Práticos):} Focada em domínio de ES6 Modules, manipulação do DOM sem bibliotecas pesadas, formulários e responsividade.
\end{itemize}

\subsection{Camada de Serviços e I/O (\texttt{services/})}

A persistência no navegador é completamente intermediada por \path{storageService.js}. Caso o navegador esteja em modo anônimo estrito ou com cotas esgotadas, o serviço captura as exceções e mantém a operação estável em memória:

\begin{boxCodigoJS}[js/services/storageService.js]
import { STORAGE_KEY_DEMANDAS, DEMANDAS_INICIAIS } from "../config/constants.js";

export function carregarDemandas() {
  try {
    const raw = localStorage.getItem(STORAGE_KEY_DEMANDAS);
    if (!raw) return DEMANDAS_INICIAIS;
    const lista = JSON.parse(raw);
    return Array.isArray(lista) && lista.length > 0 ? lista : DEMANDAS_INICIAIS;
  } catch (erro) {
    console.warn("Falha ao acessar LocalStorage; utilizando dados de contingência:", erro);
    return DEMANDAS_INICIAIS;
  }
}

export function salvarDemandas(demandas) {
  try {
    localStorage.setItem(STORAGE_KEY_DEMANDAS, JSON.stringify(demandas));
    return true;
  } catch (erro) {
    console.error("Erro crítico ao gravar demandas no LocalStorage:", erro);
    return false;
  }
}
\end{boxCodigoJS}

A comunicação externa com a planilha Google Sheets via Google Apps Script é centralizada em \path{apiService.js}. O método utiliza cabeçalho \texttt{text/plain;charset=utf-8} para evitar bloqueios de pré-voo CORS (\textit{CORS preflight}):

\begin{boxCodigoJS}[js/services/apiService.js]
export async function chamarApiAppsScript(acao, payload = {}) {
  const url = window.PROGWEB2_CONFIG?.appsScriptUrl;
  if (!url) throw new Error("URL da API do Google Apps Script não configurada.");

  const corpoRequisicao = JSON.stringify({ acao, ...payload, timestamp: new Date().toISOString() });

  const resposta = await fetch(url, {
    method: "POST",
    headers: { "Content-Type": "text/plain;charset=utf-8" },
    body: corpoRequisicao
  });

  if (!resposta.ok) throw new Error(`Falha na requisição HTTP: status ${resposta.status}`);
  return await resposta.json();
}
\end{boxCodigoJS}

\subsection{A Ponte de Retrocompatibilidade no \texttt{js/app.js}}

Uma das maiores inovações arquiteturais desta refatoração foi garantir que o \path{index.html} continuasse funcionando sem a necessidade de reescrever centenas de atributos de manipuladores declarativos legados (\texttt{onclick="..."} e \texttt{onchange="..."}).

\begin{conceito}{Escopo de Módulos vs. Escopo Global (\texttt{window})}
Ao carregar um script com \texttt{type="module"}, o navegador executa o código em escopo fechado. As funções definidas dentro do módulo \textbf{não existem no objeto global \texttt{window}}. Para viabilizar a transição sem quebras no HTML existente, o \path{js/app.js} expõe explicitamente as 79 funções manipuladoras ao escopo global.
\end{conceito}

\begin{boxCodigoJS}[js/app.js]
// Vinculação controlada das rotinas públicas ao objeto global window
const manipuladoresGlobais = {
  showTab, toggleTheme, alternarDestaqueCard,
  filtrarDemandas, abrirModalNovaDemanda, salvarDemandaFormulario,
  navegarParaEtapaSNCT, abrirModalInscricaoSNCT, confirmarInscricaoSNCT,
  abrirModalAvaliacaoRubrica, salvarAvaliacaoRubrica,
  iniciarLoginGoogle, efetuarLogoutGoogle, sincronizarComBackendAppsScript
  // ... mapeamento das 79 rotinas públicas necessárias
};

Object.entries(manipuladoresGlobais).forEach(([nome, funcao]) => {
  if (typeof funcao === "function") {
    window[nome] = funcao;
  }
});
\end{boxCodigoJS}

% =====================================================
% 4. GUIA DE MANUTENÇÃO E IMPACTO
% =====================================================
\section{Guia de Manutenção \& Métricas de Impacto}\label{sec:extensao}

\subsection{Como Adicionar uma Nova Funcionalidade ao Portal}

Para manter a integridade da arquitetura em camadas nas próximas manutenções, siga rigorosamente o fluxo de desenvolvimento abaixo:

\begin{enumerate}
  \item \textbf{Novas Constantes ou Roster:} Caso surjam novas turmas, novas chaves de cache ou novos instrumentos avaliativos, registre-os em \path{js/config/constants.js} ou \path{js/config/rubricasData.js};
  \item \textbf{Acesso a Dados:} Qualquer necessidade de consulta ou escrita em \texttt{localStorage} ou chamadas remotas de rede deve receber um método correspondente em \path{js/services/};
  \item \textbf{Renderização de Componentes:} A construção de strings HTML de cards ou tabelas pertence exclusivamente aos módulos da pasta \path{js/views/}. Use sempre \texttt{escapeHtml()} para qualquer dado digitado por usuário ou estudante;
  \item \textbf{Ponto de Conexão no DOM:} Se o novo componente demandar interação direta via atributo HTML (\texttt{onclick}), inclua a função correspondente no mapa de vinculação em \path{js/app.js}.
\end{enumerate}

\begin{atencao}{Regra de Ouro da Camada de Visualização}
Nenhum módulo da pasta \path{js/views/} deve executar diretamente comandos como \texttt{localStorage.setItem()} ou invocar \texttt{fetch()} para serviços externos. Toda a persistência deve obrigatoriamente transitar por \path{js/services/storageService.js} ou \path{js/services/apiService.js}.
\end{atencao}

\subsection{Matriz Comparativa de Métricas}

A tabela a seguir demonstra quantitativamente o ganho obtido com a refatoração modular:

\begin{table}[H]
\centering
\small
\renewcommand{\arraystretch}{1.45}\rowcolors{2}{fundo}{white}
\begin{tabularx}{\textwidth}{RYY}
\cabecalho \cab{Métrica de Qualidade} & \cab{Antes da Refatoração (Monólito)} & \cab{Após a Refatoração (ES6 Modules)} \\
Linhas totais no \texttt{index.html} & 11.070 linhas & 5.622 linhas (Redução de 49,2\%) \\
Linhas de JavaScript inline no HTML & 5.445 linhas em bloco único & 0 linhas (100\% desacoplado em módulos) \\
Organização do código & Escopo global desprotegido & 11 arquivos coesos distribuídos em 4 camadas \\
Localização e isolamento de bugs & Varredura longa e propensa a erros & Diagnóstico pontual no módulo responsável \\
Aderência pedagógica & Divergente da ementa lecionada & Idêntica à arquitetura ensinada no Roteiro 04 \\
Testabilidade de rotinas & Impossível isolar regras de negócio & Funções puras exportáveis e testáveis individualmente \\
\end{tabularx}
\caption{Comparativo das métricas de engenharia de software antes e após a modularização.}
\end{table}

\section{Checklist de Verificação Arquitetural}\label{sec:checklist}

\begin{checklist}[Checklist para validação de novas contribuições no portal]
\begin{itemize}
  \item[\caixaMarcar] \textbf{Zero scripts inline no HTML:} Nenhuma nova tag \texttt{<script>} com lógica de negócio foi adicionada ao \path{index.html}.
  \item[\caixaMarcar] \textbf{Desacoplamento do storage:} Nenhuma view executa leitura ou escrita direta no \texttt{localStorage}.
  \item[\caixaMarcar] \textbf{Sanitização contra XSS:} Todo dado dinâmico renderizado em template literals utiliza a função \texttt{escapeHtml()}.
  \item[\caixaMarcar] \textbf{Importações explícitas:} Todas as dependências entre arquivos utilizam caminhos relativos com extensão \path{.js}.
  \item[\caixaMarcar] \textbf{Ponte global registrada:} Funções chamadas por manipuladores declarativos no HTML foram exportadas em \path{js/app.js}.
  \item[\caixaMarcar] \textbf{Validação no GitHub Pages:} O portal carrega sem erros de sintaxe ou de CORS no console do navegador em produção.
\end{itemize}
\end{checklist}

\vspace{1.0em}
\begin{dica}{Referência nos Roteiros da Disciplina}
Para entender os fundamentos de importação e exportação de módulos, consulte o material de apoio em \path{roteiro-4/} (\textit{Arquitetura Modular com ES6 Modules}), e para os detalhes da integração com APIs remotas, consulte \path{roteiro-5/} (\textit{Fetch API \& Flask}).
\end{dica}

\end{document}
